/-
This file was generated by Aristotle (https://aristotle.harmonic.fun).

Lean version: leanprover/lean4:v4.24.0
Mathlib version: f897ebcf72cd16f89ab4577d0c826cd14afaafc7
This project request had uuid: 5cac3002-eea6-4251-b9bd-0838e0c9273b

To cite Aristotle, tag @Aristotle-Harmonic on GitHub PRs/issues, and add as co-author to commits:
Co-authored-by: Aristotle (Harmonic) <aristotle-harmonic@harmonic.fun>
-/

/-
We prove that the sum $\sum_{k=1}^\infty\frac{H_{2k}-H_{k}}{k^2\binom{2k}k}$ is equal to $2 \zeta(3)-\frac{\pi  \sqrt{3}\, \psi_1 \left(\frac{1}{3}\right)}{24}+\frac{\pi  \sqrt{3}\, \psi_1 \left(\frac{5}{6}\right)}{72}$, where $\psi_1$ is the trigamma function.
-/

/-
This file was generated by Aristotle (https://aristotle.harmonic.fun).

Lean version: leanprover/lean4:v4.24.0
Mathlib version: f897ebcf72cd16f89ab4577d0c826cd14afaafc7
This project request had uuid: c36a39bb-2108-4f2d-8eda-113e36c138c1

To cite Aristotle, tag @Aristotle-Harmonic on GitHub PRs/issues, and add as co-author to commits:
Co-authored-by: Aristotle (Harmonic) <aristotle-harmonic@harmonic.fun>

Aristotle's budget for this request has been reached.
If there is partial progress, it will appear in this file.
If you would like to continue working off of this partial progress,
please submit the same prompt, and add this .lean file in "Optional: Attach a Lean file as context" field.
-/

/-
This file was generated by Aristotle (https://aristotle.harmonic.fun).

Lean version: leanprover/lean4:v4.24.0
Mathlib version: f897ebcf72cd16f89ab4577d0c826cd14afaafc7
This project request had uuid: be359a99-c0c7-43b0-afbf-54c0c5c5d419

To cite Aristotle, tag @Aristotle-Harmonic on GitHub PRs/issues, and add as co-author to commits:
Co-authored-by: Aristotle (Harmonic) <aristotle-harmonic@harmonic.fun>

Aristotle's budget for this request has been reached.
If there is partial progress, it will appear in this file.
If you would like to continue working off of this partial progress,
please submit the same prompt, and add this .lean file in "Optional: Attach a Lean file as context" field.
-/

import Mathlib

set_option linter.mathlibStandardSet false

open scoped BigOperators
open scoped Real
open scoped Nat
open scoped Classical
open scoped Pointwise

set_option maxHeartbeats 0
set_option maxRecDepth 4000
set_option synthInstance.maxHeartbeats 20000
set_option synthInstance.maxSize 128

set_option relaxedAutoImplicit false
set_option autoImplicit false

noncomputable section

#check harmonic
#check riemannZeta
#check Real.pi
#check Real.sqrt

#check ∑' (n : ℕ), 1 / (n + 1 : ℝ)^2

/-
Definition of the trigamma function as a series.
-/
noncomputable def trigamma (z : ℝ) : ℝ := ∑' (n : ℕ), 1 / (n + z)^2

#check ∫ x in (0 : ℝ)..(1 : ℝ), x

/-
Definitions of the function F(z) and the sum S.
-/
noncomputable def F (z : ℝ) : ℝ := ∑' k : ℕ, if k = 0 then 0 else z^k / ((k : ℝ)^2 * (Nat.choose (2 * k) k : ℝ))

noncomputable def S_sum : ℝ := ∑' k : ℕ, if k = 0 then 0 else (harmonic (2 * k) - harmonic k : ℝ) / ((k : ℝ) ^ 2 * (Nat.choose (2 * k) k : ℝ))

/-
The inverse of the central binomial coefficient binom(2k, k) is equal to k times the integral from 0 to 1 of t^(k-1) * (1-t)^k.
-/
lemma inv_binom_eq_integral (k : ℕ) (hk : k ≠ 0) :
  (1 : ℝ) / (Nat.choose (2 * k) k) = k * ∫ t in (0 : ℝ)..1, t^(k - 1) * (1 - t)^k := by
    -- Use the Beta function identity: $\int_0^1 t^{k-1}(1-t)^k \, dt = \frac{\Gamma(k)\Gamma(k+1)}{\Gamma(2k+1)}$.
    have h_beta_identity : ∫ (t : ℝ) in (0 : ℝ)..1, t ^ (k - 1) * (1 - t) ^ k = (Nat.factorial (k - 1) * Nat.factorial k : ℝ) / Nat.factorial (2 * k) := by
      have h_beta : ∀ a b : ℕ, ∫ x in (0 : ℝ)..1, x ^ a * (1 - x) ^ b = (Nat.factorial a * Nat.factorial b : ℝ) / Nat.factorial (a + b + 1) := by
        intro a b; induction' b with b ih generalizing a <;> norm_num [ Nat.factorial ] at *;
        · rw [ inv_eq_one_div, div_eq_div_iff ] <;> ring <;> positivity;
        · -- Use integration by parts.
          have h_parts : ∀ a b : ℕ, ∫ x in (0 : ℝ)..1, x ^ a * (1 - x) ^ (b + 1) = (∫ x in (0 : ℝ)..1, x ^ a * (1 - x) ^ b) - (∫ x in (0 : ℝ)..1, x ^ (a + 1) * (1 - x) ^ b) := by
            intro a b; rw [ ← intervalIntegral.integral_sub ( by exact Continuous.intervalIntegrable ( by continuity ) _ _ ) ( by exact Continuous.intervalIntegrable ( by continuity ) _ _ ) ] ; congr ; ext ; ring;
          rw [ h_parts, ih, ih ];
          field_simp;
          simpa [ Nat.add_comm, Nat.add_left_comm, Nat.factorial ] using by ring;
      convert h_beta ( k - 1 ) k using 2 ; cases k <;> norm_num [ Nat.succ_mul ] ; ring;
    rw [ h_beta_identity, Nat.cast_choose ];
    · cases k <;> norm_num [ two_mul, Nat.factorial ] at * ; ring;
    · linarith

/-
F(z) is equal to the integral from 0 to 1 of -log(1 - z*t*(1-t))/t.
-/
lemma F_eq_integral (z : ℝ) (hz : z ∈ Set.Icc 0 1) : F z = ∫ t in (0 : ℝ)..1, - Real.log (1 - z * t * (1 - t)) / t := by
  -- Using the identity $\frac{1}{\binom{2k}{k}} = k \int_0^1 t^{k-1}(1-t)^k dt$, we get
  have h_identity : ∀ k : ℕ, k ≠ 0 → (z^k / ((k : ℝ)^2 * (Nat.choose (2 * k) k : ℝ))) = ∫ t in (0 : ℝ)..1, z^k * t^(k - 1) * (1 - t)^k / k := by
    intro k hk; have := inv_binom_eq_integral k hk; simp_all +decide [ div_eq_mul_inv, mul_assoc, mul_comm, mul_left_comm, pow_succ, tsum_mul_left ] ;
  -- By Fubini's theorem, we can interchange the sum and the integral.
  have h_fubini : ∫ t in (0 : ℝ)..1, ∑' k : ℕ, (if k = 0 then 0 else z^k * t^(k - 1) * (1 - t)^k / k) = ∑' k : ℕ, ∫ t in (0 : ℝ)..1, (if k = 0 then 0 else z^k * t^(k - 1) * (1 - t)^k / k) := by
    rw [ intervalIntegral.integral_of_le zero_le_one, MeasureTheory.integral_tsum ];
    · norm_num [ intervalIntegral.integral_of_le zero_le_one ];
    · exact fun k => Measurable.aestronglyMeasurable ( by split_ifs <;> [ exact measurable_const; exact Measurable.mul ( Measurable.mul ( measurable_const.mul ( measurable_id.pow_const _ ) ) ( measurable_const.sub measurable_id |> Measurable.pow_const <| _ ) ) measurable_const ] );
    · refine' ne_of_lt ( lt_of_le_of_lt ( ENNReal.tsum_le_tsum _ ) _ );
      use fun k => ENNReal.ofReal ( ∫ t in Set.Ioc 0 1, z ^ k * t ^ ( k - 1 ) * ( 1 - t ) ^ k / k );
      · intro k; by_cases hk : k = 0 <;> simp_all +decide [ MeasureTheory.ofReal_integral_eq_lintegral_ofReal ] ;
        rw [ MeasureTheory.ofReal_integral_eq_lintegral_ofReal ];
        · refine' MeasureTheory.setLIntegral_mono' measurableSet_Ioc fun x hx => _;
          rw [ Real.enorm_eq_ofReal ( div_nonneg ( mul_nonneg ( mul_nonneg ( pow_nonneg hz.1 _ ) ( pow_nonneg hx.1.le _ ) ) ( pow_nonneg ( sub_nonneg.2 hx.2 ) _ ) ) ( Nat.cast_nonneg _ ) ) ];
        · exact Continuous.integrableOn_Ioc ( by continuity );
        · filter_upwards [ MeasureTheory.ae_restrict_mem measurableSet_Ioc ] with t ht using div_nonneg ( mul_nonneg ( mul_nonneg ( pow_nonneg hz.1 _ ) ( pow_nonneg ht.1.le _ ) ) ( pow_nonneg ( sub_nonneg.2 ht.2 ) _ ) ) ( Nat.cast_nonneg _ );
      · -- We'll use the fact that the series converges uniformly to interchange the sum and integral.
        have h_uniform : Summable (fun k : ℕ => ∫ t in Set.Ioc (0 : ℝ) 1, z ^ k * t ^ (k - 1) * (1 - t) ^ k / k) := by
          have h_uniform : Summable (fun k : ℕ => z ^ k / ((k : ℝ) ^ 2 * (Nat.choose (2 * k) k : ℝ))) := by
            refine' Summable.of_nonneg_of_le ( fun k => div_nonneg ( pow_nonneg hz.1 _ ) ( mul_nonneg ( sq_nonneg _ ) ( Nat.cast_nonneg _ ) ) ) ( fun k => _ ) ( Real.summable_one_div_nat_pow.2 one_lt_two );
            rcases eq_or_ne k 0 <;> simp_all +decide [ division_def ];
            have := h_identity k ‹_›; rw [ ← this ] ; ring_nf; norm_num [ ‹k ≠ 0› ] ;
            exact mul_le_of_le_one_left ( by positivity ) ( mul_le_one₀ ( pow_le_one₀ hz.1 hz.2 ) ( inv_nonneg.2 ( Nat.cast_nonneg _ ) ) ( inv_le_one_of_one_le₀ ( mod_cast Nat.choose_pos ( by linarith ) ) ) );
          exact h_uniform.congr fun k => by by_cases hk : k = 0 <;> simp_all +decide [ intervalIntegral.integral_of_le zero_le_one ] ;
        exact?;
  convert h_fubini.symm using 1;
  · exact tsum_congr fun k => by aesop;
  · -- By Fubini's theorem, we can interchange the sum and the integral in the expression.
    have h_fubini : ∀ t ∈ Set.Ioo (0 : ℝ) 1, -Real.log (1 - z * t * (1 - t)) / t = ∑' k : ℕ, (if k = 0 then 0 else z^k * t^(k - 1) * (1 - t)^k / k) := by
      intro t ht
      have h_series : -Real.log (1 - z * t * (1 - t)) = ∑' k : ℕ, (z^k * t^k * (1 - t)^k) / k := by
        have := Real.hasSum_pow_div_log_of_abs_lt_one ( show |z * t * ( 1 - t )| < 1 from ?_ );
        · rw [ Summable.tsum_eq_zero_add ];
          · simpa [ mul_pow ] using this.tsum_eq.symm;
          · exact summable_nat_add_iff 1 |>.1 <| by simpa [ mul_pow ] using this.summable;
        · exact abs_lt.mpr ⟨ by nlinarith [ hz.1, hz.2, ht.1, ht.2, mul_nonneg hz.1 ht.1.le ], by nlinarith [ hz.1, hz.2, ht.1, ht.2, mul_nonneg hz.1 ht.1.le, mul_nonneg hz.1 ( sub_nonneg.mpr ht.2.le ) ] ⟩;
      rw [ h_series, div_eq_iff ht.1.ne' ];
      rw [ ← tsum_mul_right ] ; congr ; ext k ; cases k <;> norm_num [ pow_succ, mul_assoc, mul_comm, mul_left_comm ] ; ring;
    rw [ intervalIntegral.integral_of_le zero_le_one, intervalIntegral.integral_of_le zero_le_one ];
    simpa only [ MeasureTheory.integral_Ioc_eq_integral_Ioo ] using MeasureTheory.setIntegral_congr_fun measurableSet_Ioo h_fubini

/-
S is equal to the integral from 0 to 1 of (F(x) - F(x^2)) / (1 - x).
-/
lemma S_eq_integral : S_sum = ∫ x in (0 : ℝ)..1, (F x - F (x^2)) / (1 - x) := by
  -- By Fubini's theorem, we can interchange the sum and the integral.
  have h_fubini : ∫ x in (0 : ℝ)..1, (F x - F (x ^ 2)) / (1 - x) = ∑' k : ℕ, (if k = 0 then 0 else (1 / ((k : ℝ) ^ 2 * (Nat.choose (2 * k) k : ℝ))) * (∫ x in (0 : ℝ)..1, (x ^ k - x ^ (2 * k)) / (1 - x))) := by
    -- By Fubini's theorem, we can interchange the sum and integral.
    have h_fubini : ∫ x in (0 : ℝ)..1, ∑' k : ℕ, (if k = 0 then 0 else (1 / ((k : ℝ) ^ 2 * (Nat.choose (2 * k) k : ℝ))) * ((x ^ k - x ^ (2 * k)) / (1 - x))) = ∑' k : ℕ, (if k = 0 then 0 else (1 / ((k : ℝ) ^ 2 * (Nat.choose (2 * k) k : ℝ))) * (∫ x in (0 : ℝ)..1, (x ^ k - x ^ (2 * k)) / (1 - x))) := by
      rw [ intervalIntegral.integral_of_le zero_le_one, MeasureTheory.integral_tsum ];
      · norm_num [ intervalIntegral.integral_of_le zero_le_one, MeasureTheory.integral_const_mul ];
        exact tsum_congr fun n => by split_ifs <;> simp +decide [ *, MeasureTheory.integral_const_mul ] ;
      · exact fun i => Measurable.aestronglyMeasurable ( by exact Measurable.ite ( by norm_num ) measurable_const <| by exact Measurable.mul ( measurable_const ) <| by exact Measurable.div ( by exact Measurable.sub ( measurable_id.pow_const _ ) <| measurable_id.pow_const _ ) <| by exact measurable_const.sub measurable_id );
      · -- We'll use the fact that if the series converges absolutely, then the integral of the absolute value also converges.
        have h_abs_conv : Summable (fun k : ℕ => (1 : ℝ) / ((k : ℝ) ^ 2 * (Nat.choose (2 * k) k : ℝ)) * ∫ x in (0 : ℝ)..1, |(x ^ k - x ^ (2 * k)) / (1 - x)|) := by
          -- We'll use the fact that |(x^k - x^(2k)) / (1 - x)| ≤ k for x in [0, 1].
          have h_bound : ∀ k : ℕ, k ≠ 0 → ∫ x in (0 : ℝ)..1, |(x ^ k - x ^ (2 * k)) / (1 - x)| ≤ k := by
            -- We'll use the fact that |(x^k - x^(2k)) / (1 - x)| ≤ k for x in [0, 1]. This follows from the fact that the integrand is bounded above by k.
            intros k hk_nonzero
            have h_bound : ∀ x ∈ Set.Ioo (0 : ℝ) 1, |(x ^ k - x ^ (2 * k)) / (1 - x)| ≤ k := by
              -- We can factor the numerator as $x^k(1 - x^k)$ and simplify the expression.
              intro x hx
              have h_factor : (x ^ k - x ^ (2 * k)) / (1 - x) = x ^ k * (∑ i ∈ Finset.range k, x ^ i) := by
                rw [ div_eq_iff ] <;> ring;
                · rw [ pow_mul ] ; nlinarith [ geom_sum_mul x k, hx.1, hx.2, pow_pos hx.1 k ] ;
                · linarith [ hx.1, hx.2 ];
              rw [ h_factor, abs_of_nonneg ( mul_nonneg ( pow_nonneg hx.1.le _ ) ( Finset.sum_nonneg fun _ _ => pow_nonneg hx.1.le _ ) ) ];
              exact le_trans ( mul_le_of_le_one_left ( Finset.sum_nonneg fun _ _ => pow_nonneg hx.1.le _ ) ( pow_le_one₀ hx.1.le hx.2.le ) ) ( le_trans ( Finset.sum_le_sum fun _ _ => pow_le_one₀ hx.1.le hx.2.le ) ( by norm_num ) );
            rw [ intervalIntegral.integral_of_le zero_le_one, MeasureTheory.integral_Ioc_eq_integral_Ioo ];
            refine' le_trans ( MeasureTheory.integral_mono_of_nonneg _ _ _ ) _;
            exacts [ fun _ => k, Filter.Eventually.of_forall fun _ => abs_nonneg _, Continuous.integrableOn_Icc ( by continuity ) |> fun h => h.mono_set <| Set.Ioo_subset_Icc_self, Filter.eventually_of_mem ( MeasureTheory.ae_restrict_mem measurableSet_Ioo ) fun x hx => h_bound x hx, by norm_num ];
          -- Using the bound, we can show that the series converges absolutely.
          have h_abs_conv : Summable (fun k : ℕ => (1 : ℝ) / ((k : ℝ) ^ 2 * (Nat.choose (2 * k) k : ℝ)) * k) := by
            -- We'll use the fact that $\binom{2k}{k} \geq 2^k$ for all $k$.
            have h_binom : ∀ k : ℕ, (Nat.choose (2 * k) k : ℝ) ≥ 2 ^ k := by
              intro k; norm_cast; induction' k with k ih <;> norm_num [ Nat.pow_succ', Nat.mul_succ, Nat.choose_succ_succ ] at *;
              linarith [ Nat.choose_le_succ ( 2 * k ) k ];
            -- Using the bound, we can show that the series converges absolutely by comparing it to a convergent p-series.
            have h_abs_conv : Summable (fun k : ℕ => (1 : ℝ) / ((k : ℝ) * (2 ^ k))) := by
              norm_num +zetaDelta at *;
              exact Summable.of_nonneg_of_le ( fun k => by positivity ) ( fun k => mul_le_of_le_one_right ( by positivity ) ( by cases k <;> simpa using inv_le_one_of_one_le₀ <| mod_cast Nat.le_add_left _ _ ) ) <| by simpa using summable_geometric_two;
            refine Summable.of_nonneg_of_le ( fun k => by positivity ) ( fun k => ?_ ) h_abs_conv ; by_cases hk : k = 0 <;> simp_all +decide [ sq, mul_assoc, mul_comm, mul_left_comm ];
            exact mul_le_mul_of_nonneg_left ( inv_anti₀ ( by positivity ) ( h_binom k ) ) ( by positivity );
          refine' .of_nonneg_of_le ( fun k => mul_nonneg ( by positivity ) ( intervalIntegral.integral_nonneg ( by norm_num ) fun x hx => abs_nonneg _ ) ) ( fun k => mul_le_mul_of_nonneg_left ( if hk : k = 0 then by norm_num [ hk ] else h_bound k hk ) ( by positivity ) ) h_abs_conv;
        refine' ne_of_lt ( lt_of_le_of_lt ( ENNReal.tsum_le_tsum fun k => _ ) _ );
        use fun k => ENNReal.ofReal ( ∫ x in Set.Ioc 0 1, |(if k = 0 then 0 else 1 / ((k : ℝ) ^ 2 * (Nat.choose (2 * k) k : ℝ)) * ((x ^ k - x ^ (2 * k)) / (1 - x)))| );
        · rw [ MeasureTheory.ofReal_integral_eq_lintegral_ofReal ];
          · refine' MeasureTheory.lintegral_mono fun x => _;
            rw [ ENNReal.le_ofReal_iff_toReal_le ] <;> norm_num;
          · by_cases hk : k = 0 <;> simp +decide [ hk ];
            refine' MeasureTheory.Integrable.const_mul _ _;
            refine' MeasureTheory.Integrable.abs _;
            -- The integral of $(x^k - x^{2k}) / (1 - x)$ over $(0, 1)$ is convergent because it is a rational function with a simple pole at $x = 1$.
            have h_integrable : MeasureTheory.IntegrableOn (fun x : ℝ => (x ^ k - x ^ (2 * k)) / (1 - x)) (Set.Ioo 0 1) := by
              have h_factor : ∀ x : ℝ, x ∈ Set.Ioo 0 1 → (x ^ k - x ^ (2 * k)) / (1 - x) = x ^ k * (∑ i ∈ Finset.range k, x ^ i) := by
                intro x hx; rw [ div_eq_iff ( by linarith [ hx.1, hx.2 ] ) ] ; ring;
                rw [ pow_mul ] ; nlinarith [ hx.1, hx.2, geom_sum_mul x k, pow_pos hx.1 k ] ;
              rw [ MeasureTheory.integrableOn_congr_fun ( fun x hx => h_factor x hx ) measurableSet_Ioo ];
              exact Continuous.integrableOn_Icc ( by continuity ) |> fun h => h.mono_set <| Set.Ioo_subset_Icc_self;
            rwa [ MeasureTheory.IntegrableOn, MeasureTheory.Measure.restrict_congr_set MeasureTheory.Ioo_ae_eq_Ioc ] at *;
          · exact Filter.Eventually.of_forall fun x => abs_nonneg _;
        · rw [ ← ENNReal.ofReal_tsum_of_nonneg ] <;> norm_num;
          · exact fun n => MeasureTheory.integral_nonneg fun x => abs_nonneg _;
          · refine' .of_nonneg_of_le ( fun k => MeasureTheory.integral_nonneg fun x => abs_nonneg _ ) ( fun k => _ ) h_abs_conv;
            by_cases hk : k = 0 <;> simp +decide [ hk, mul_assoc, mul_comm, mul_left_comm, ← intervalIntegral.integral_of_le zero_le_one ];
    convert h_fubini using 1;
    refine' intervalIntegral.integral_congr fun x hx => _;
    rw [ show F x = ∑' k : ℕ, if k = 0 then 0 else x ^ k / ( ( k : ℝ ) ^ 2 * Nat.choose ( 2 * k ) k ) from ?_, show F ( x ^ 2 ) = ∑' k : ℕ, if k = 0 then 0 else x ^ ( 2 * k ) / ( ( k : ℝ ) ^ 2 * Nat.choose ( 2 * k ) k ) from ?_ ];
    · rw [ ← Summable.tsum_sub ];
      · rw [ ← tsum_div_const ] ; congr ; ext k ; split_ifs <;> simp +decide [ *, div_eq_mul_inv, mul_assoc, mul_comm, mul_left_comm ];
        ring;
      · refine' Summable.of_nonneg_of_le ( fun k => _ ) ( fun k => _ ) ( Real.summable_one_div_nat_pow.2 one_lt_two );
        · split_ifs <;> first | positivity | exact div_nonneg ( pow_nonneg ( by aesop ) _ ) ( by positivity );
        · split_ifs <;> norm_num;
          field_simp;
          exact div_le_one_of_le₀ ( le_trans ( pow_le_one₀ ( by aesop ) ( by aesop ) ) ( mod_cast Nat.choose_pos ( by linarith ) ) ) ( Nat.cast_nonneg _ );
      · refine' Summable.of_nonneg_of_le ( fun k => _ ) ( fun k => _ ) ( Real.summable_one_div_nat_pow.2 one_lt_two );
        · split_ifs <;> first | positivity | exact div_nonneg ( pow_nonneg ( by aesop ) _ ) ( by positivity ) ;
        · split_ifs <;> norm_num at *;
          exact le_trans ( mul_le_of_le_one_left ( by positivity ) ( pow_le_one₀ hx.1 hx.2 ) ) ( inv_anti₀ ( by positivity ) ( le_mul_of_one_le_right ( by positivity ) ( mod_cast Nat.choose_pos ( by linarith ) ) ) );
    · exact tsum_congr fun k => by split_ifs <;> ring;
    · exact?;
  -- We'll use the fact that $\int_0^1 \frac{x^k - x^{2k}}{1-x} \, dx = H_{2k} - H_k$.
  have h_integral : ∀ k : ℕ, k ≠ 0 → ∫ x in (0 : ℝ)..1, (x ^ k - x ^ (2 * k)) / (1 - x) = (harmonic (2 * k) - harmonic k : ℝ) := by
    -- We'll use the fact that $\int_0^1 \frac{x^k - x^{2k}}{1-x} \, dx = \int_0^1 \sum_{n=k}^{2k-1} x^n \, dx$.
    intro k hk_ne_zero
    have h_integral_sum : ∫ x in (0 : ℝ)..1, (x ^ k - x ^ (2 * k)) / (1 - x) = ∫ x in (0 : ℝ)..1, ∑ n ∈ Finset.Ico k (2 * k), x ^ n := by
      have h_integral_sum : ∀ x ∈ Set.Ioo (0 : ℝ) 1, (x ^ k - x ^ (2 * k)) / (1 - x) = ∑ n ∈ Finset.Ico k (2 * k), x ^ n := by
        intro x hx; rw [ div_eq_iff ( by linarith [ hx.1, hx.2 ] ) ] ; rw [ Finset.sum_Ico_eq_sub _ ( by linarith [ Nat.pos_of_ne_zero hk_ne_zero ] ) ] ; ring;
        have := geom_sum_mul_neg x ( k * 2 ) ; have := geom_sum_mul_neg x k ; norm_num [ pow_mul ] at * ; nlinarith [ hx.1, hx.2 ] ;
      rw [ intervalIntegral.integral_of_le zero_le_one, intervalIntegral.integral_of_le zero_le_one ] ; rw [ MeasureTheory.integral_Ioc_eq_integral_Ioo, MeasureTheory.integral_Ioc_eq_integral_Ioo ] ; exact MeasureTheory.setIntegral_congr_fun measurableSet_Ioo h_integral_sum;
    rw [ h_integral_sum, intervalIntegral.integral_finset_sum ] <;> norm_num [ hk_ne_zero ];
    norm_num [ two_mul, Finset.sum_Ico_eq_sub, harmonic ];
  rw [ h_fubini, tsum_congr ] ; aesop;
  grind

/-
The derivative of F(z) is equal to the integral from 0 to 1 of (1-t)/(1 - z*t*(1-t)).
-/
lemma F_deriv_eq_integral (z : ℝ) (hz : z ∈ Set.Ioo 0 1) :
  deriv F z = ∫ t in (0 : ℝ)..1, (1 - t) / (1 - z * t * (1 - t)) := by
    refine' HasDerivAt.deriv _;
    -- We'll use the fact that if the series of derivatives converges uniformly, then the derivative of the sum is the sum of the derivatives.
    have h_deriv_series : Filter.Tendsto (fun h => (∫ t in (0 : ℝ)..1, -Real.log (1 - (z + h) * t * (1 - t)) / t - (-Real.log (1 - z * t * (1 - t)) / t)) / h) (nhdsWithin 0 {0}ᶜ) (nhds (∫ t in (0 : ℝ)..1, (1 - t) / (1 - z * t * (1 - t)))) := by
      -- We'll use the fact that if the series of derivatives converges uniformly, then the derivative of the sum is the sum of the derivatives. We'll show that the series of derivatives converges uniformly.
      have h_uniform_converges : Filter.Tendsto (fun h => ∫ t in (0 : ℝ)..1, (-(Real.log (1 - (z + h) * t * (1 - t))) / t - (-(Real.log (1 - z * t * (1 - t))) / t)) / h) (nhdsWithin 0 {0}ᶜ) (nhds (∫ t in (0 : ℝ)..1, (1 - t) / (1 - z * t * (1 - t)))) := by
        -- To apply the Dominated Convergence Theorem, we need to show that the sequence of functions is dominated by an integrable function.
        have h_dominate : ∀ h : ℝ, |h| < (1 - z) / 2 → ∀ t ∈ Set.Ioo (0 : ℝ) 1, |(-Real.log (1 - (z + h) * t * (1 - t)) / t - -Real.log (1 - z * t * (1 - t)) / t) / h| ≤ 2 := by
          intros h hh t ht
          have h_abs : |(-Real.log (1 - (z + h) * t * (1 - t)) / t - -Real.log (1 - z * t * (1 - t)) / t) / h| ≤ 2 := by
            have h_mean_value : ∃ c ∈ Set.Icc (min z (z + h)) (max z (z + h)), (-Real.log (1 - (z + h) * t * (1 - t)) / t - -Real.log (1 - z * t * (1 - t)) / t) = (1 - t) / (1 - c * t * (1 - t)) * h := by
              cases eq_or_ne h 0 <;> simp_all +decide [ div_eq_inv_mul ];
              -- Apply the mean value theorem to the interval [z, z + h].
              obtain ⟨c, hc⟩ : ∃ c ∈ Set.Ioo (min z (z + h)) (max z (z + h)), deriv (fun x => -Real.log (1 - x * t * (1 - t)) / t) c = ( -Real.log (1 - (z + h) * t * (1 - t)) / t - -Real.log (1 - z * t * (1 - t)) / t ) / h := by
                cases max_cases z ( z + h ) <;> cases min_cases z ( z + h ) <;> simp_all +decide [ div_eq_inv_mul ];
                · have := exists_deriv_eq_slope ( f := fun x => -t⁻¹ * Real.log ( 1 - x * t * ( 1 - t ) ) ) ( show z + h < z by linarith );
                  norm_num +zetaDelta at *;
                  exact this ( ContinuousOn.neg <| ContinuousOn.mul continuousOn_const <| ContinuousOn.log ( ContinuousOn.sub continuousOn_const <| ContinuousOn.mul ( continuousOn_id.mul continuousOn_const ) <| continuousOn_const.sub continuousOn_const ) fun x hx => by nlinarith [ hx.1, hx.2, abs_lt.mp hh, mul_pos ht.1 ( sub_pos.mpr ht.2 ) ] ) ( DifferentiableOn.mul ( differentiableOn_const _ ) <| DifferentiableOn.log ( DifferentiableOn.sub ( differentiableOn_const _ ) <| DifferentiableOn.mul ( differentiableOn_id.mul_const _ ) <| differentiableOn_const _ |> DifferentiableOn.sub <| differentiableOn_const _ ) fun x hx => by nlinarith [ hx.1, hx.2, abs_lt.mp hh, mul_pos ht.1 ( sub_pos.mpr ht.2 ) ] ) |> fun ⟨ c, hc₁, hc₂ ⟩ => ⟨ c, hc₁, by rw [ hc₂ ] ; ring ⟩;
                · have := exists_deriv_eq_slope ( f := fun x => -t⁻¹ * Real.log ( 1 - x * t * ( 1 - t ) ) ) ( show z < z + h by linarith );
                  norm_num +zetaDelta at *;
                  exact this ( ContinuousOn.neg <| ContinuousOn.mul continuousOn_const <| ContinuousOn.log ( ContinuousOn.sub continuousOn_const <| ContinuousOn.mul ( continuousOn_id.mul continuousOn_const ) continuousOn_const ) fun x hx => by nlinarith [ hx.1, hx.2, abs_lt.mp hh, mul_pos ht.1 ( sub_pos.mpr ht.2 ) ] ) ( DifferentiableOn.mul ( differentiableOn_const _ ) <| DifferentiableOn.log ( DifferentiableOn.sub ( differentiableOn_const _ ) <| DifferentiableOn.mul ( differentiableOn_id.mul_const _ ) <| differentiableOn_const _ ) fun x hx => by nlinarith [ hx.1, hx.2, abs_lt.mp hh, mul_pos ht.1 ( sub_pos.mpr ht.2 ) ] ) |> fun ⟨ c, hc₁, hc₂ ⟩ => ⟨ c, hc₁, by linear_combination hc₂ ⟩;
              norm_num [ show 1 - c * t * ( 1 - t ) ≠ 0 from by cases max_cases z ( z + h ) <;> cases min_cases z ( z + h ) <;> nlinarith [ abs_lt.mp hh, hc.1.1, hc.1.2, mul_pos ht.1 ( sub_pos.mpr ht.2 ) ] ] at *;
              rw [ deriv.log ] at hc <;> norm_num at *;
              · grind;
              · cases hc.1.1 <;> cases hc.1.2 <;> nlinarith [ abs_lt.mp hh, mul_pos ht.1 ( sub_pos.mpr ht.2 ) ]
            obtain ⟨ c, hc₁, hc₂ ⟩ := h_mean_value; rw [ hc₂ ] ; by_cases hh : h = 0 <;> simp_all +decide [ abs_div, abs_mul ] ;
            rw [ div_le_iff₀ ] <;> norm_num [ abs_of_pos, ht.1, ht.2 ];
            · cases abs_cases ( 1 - c * t * ( 1 - t ) ) <;> cases hc₁.1 <;> cases hc₁.2 <;> nlinarith [ mul_pos ht.1 ( sub_pos.2 ht.2 ), abs_lt.mp ‹_› ];
            · cases hc₁.1 <;> cases hc₁.2 <;> nlinarith [ mul_pos ht.1 ( sub_pos.2 ht.2 ), abs_lt.mp ‹_› ];
          exact h_abs;
        -- Apply the Dominated Convergence Theorem.
        have h_dominated_convergence : Filter.Tendsto (fun h => ∫ t in (Set.Ioo (0 : ℝ) 1), (-Real.log (1 - (z + h) * t * (1 - t)) / t - -Real.log (1 - z * t * (1 - t)) / t) / h) (nhdsWithin 0 {0}ᶜ) (nhds (∫ t in (Set.Ioo (0 : ℝ) 1), (1 - t) / (1 - z * t * (1 - t)))) := by
          refine' MeasureTheory.tendsto_integral_filter_of_dominated_convergence _ _ _ _ _;
          refine' fun t => 2;
          · refine' Filter.eventually_of_mem self_mem_nhdsWithin fun n hn => Measurable.aestronglyMeasurable _;
            fun_prop (disch := norm_num);
          · rw [ eventually_nhdsWithin_iff ];
            filter_upwards [ Metric.ball_mem_nhds _ ( show 0 < ( 1 - z ) / 2 by linarith [ hz.1, hz.2 ] ) ] with x hx hx' using Filter.eventually_of_mem ( MeasureTheory.ae_restrict_mem measurableSet_Ioo ) fun t ht => h_dominate x ( by simpa using hx ) t ht;
          · norm_num;
          · -- We'll use the fact that the derivative of $-\frac{\log(1 - z t (1 - t))}{t}$ with respect to $z$ is $\frac{1 - t}{1 - z t (1 - t)}$.
            have h_deriv : ∀ t ∈ Set.Ioo (0 : ℝ) 1, HasDerivAt (fun z => -Real.log (1 - z * t * (1 - t)) / t) ((1 - t) / (1 - z * t * (1 - t))) z := by
              intro t ht; convert HasDerivAt.div_const ( HasDerivAt.neg ( HasDerivAt.log ( HasDerivAt.sub ( hasDerivAt_const _ _ ) ( HasDerivAt.mul ( HasDerivAt.mul ( hasDerivAt_id' z ) ( hasDerivAt_const _ _ ) ) ( hasDerivAt_const _ _ ) ) ) _ ) ) _ using 1 <;> norm_num ; ring;
              · simpa [ sq, mul_assoc, mul_comm t, ht.1.ne' ] using by ring;
              · nlinarith [ hz.1, hz.2, ht.1, ht.2, mul_pos hz.1 ht.1, mul_pos hz.1 ( sub_pos.2 ht.2 ), mul_pos ht.1 ( sub_pos.2 hz.2 ) ];
            filter_upwards [ MeasureTheory.ae_restrict_mem measurableSet_Ioo ] with t ht using by simpa [ div_eq_inv_mul ] using h_deriv t ht |> HasDerivAt.tendsto_slope_zero;
        simpa only [ MeasureTheory.integral_Ioc_eq_integral_Ioo, intervalIntegral.integral_of_le zero_le_one ] using h_dominated_convergence;
      simpa only [ intervalIntegral.integral_div ] using h_uniform_converges;
    have h_eq : ∀ h, abs h < min z (1 - z) → F (z + h) - F z = ∫ t in (0 : ℝ)..1, -Real.log (1 - (z + h) * t * (1 - t)) / t - (-Real.log (1 - z * t * (1 - t)) / t) := by
      intro h hh;
      rw [ intervalIntegral.integral_sub ];
      · rw [ F_eq_integral, F_eq_integral ];
        · exact ⟨ hz.1.le, hz.2.le ⟩;
        · constructor <;> linarith [ abs_lt.mp hh, min_le_left z ( 1 - z ), min_le_right z ( 1 - z ) ];
      · rw [ intervalIntegrable_iff_integrableOn_Ioo_of_le ];
        · -- We'll use the fact that $-\log(1 - (z + h) t (1 - t)) / t$ is bounded on $(0, 1)$.
          have h_bound : ∀ t ∈ Set.Ioo 0 1, abs (-Real.log (1 - (z + h) * t * (1 - t)) / t) ≤ 2 / (1 - (z + h)) := by
            intros t ht
            have h_log_bound : abs (Real.log (1 - (z + h) * t * (1 - t))) ≤ (z + h) * t * (1 - t) / (1 - (z + h)) := by
              have h_log_bound : ∀ x ∈ Set.Ioo 0 1, abs (Real.log (1 - x)) ≤ x / (1 - x) := by
                intro x hx; rw [ abs_of_nonpos ( Real.log_nonpos ( by linarith [ hx.1, hx.2 ] ) ( by linarith [ hx.1, hx.2 ] ) ) ] ; rw [ div_eq_mul_inv ] ; nlinarith [ hx.1, hx.2, Real.log_inv ( 1 - x ), Real.log_le_sub_one_of_pos ( inv_pos.mpr ( by linarith [ hx.1, hx.2 ] : 0 < 1 - x ) ), mul_inv_cancel₀ ( by linarith [ hx.1, hx.2 ] : ( 1 - x ) ≠ 0 ) ] ;
              refine le_trans ( h_log_bound _ ?_ ) ?_;
              · constructor <;> nlinarith [ abs_lt.mp hh, min_le_left z ( 1 - z ), min_le_right z ( 1 - z ), ht.1, ht.2, mul_pos ht.1 ( sub_pos.mpr ht.2 ) ];
              · gcongr;
                · exact mul_nonneg ( mul_nonneg ( by linarith [ abs_lt.mp hh, min_le_left z ( 1 - z ), min_le_right z ( 1 - z ) ] ) ht.1.le ) ( sub_nonneg.mpr ht.2.le );
                · linarith [ abs_lt.mp hh, min_le_left z ( 1 - z ), min_le_right z ( 1 - z ) ];
                · nlinarith [ abs_lt.mp hh, min_le_left z ( 1 - z ), min_le_right z ( 1 - z ), ht.1, ht.2, mul_pos ht.1 ( sub_pos.mpr ht.2 ), mul_pos ht.1 ( sub_pos.mpr ( show z + h > 0 from by linarith [ abs_lt.mp hh, min_le_left z ( 1 - z ), min_le_right z ( 1 - z ), hz.1, hz.2 ] ) ), mul_pos ( sub_pos.mpr ht.2 ) ( sub_pos.mpr ( show z + h > 0 from by linarith [ abs_lt.mp hh, min_le_left z ( 1 - z ), min_le_right z ( 1 - z ), hz.1, hz.2 ] ) ) ];
            rw [ abs_div, abs_neg, abs_of_nonneg ht.1.le ];
            rw [ div_le_iff₀ ht.1 ];
            refine le_trans h_log_bound ?_;
            rw [ div_mul_eq_mul_div, div_le_div_iff_of_pos_right ] <;> nlinarith [ abs_lt.mp hh, min_le_left z ( 1 - z ), min_le_right z ( 1 - z ), ht.1, ht.2, mul_pos ht.1 ( sub_pos.mpr ht.2 ) ];
          refine' MeasureTheory.Integrable.mono' _ _ _;
          refine' fun t => 2 / ( 1 - ( z + h ) );
          · norm_num;
          · exact Measurable.aestronglyMeasurable ( by exact Measurable.mul ( Measurable.neg ( Measurable.log ( by exact Measurable.sub measurable_const ( by exact Measurable.mul ( measurable_const.mul measurable_id' ) ( measurable_const.sub measurable_id' ) ) ) ) ) ( measurable_id'.inv ) );
          · filter_upwards [ MeasureTheory.ae_restrict_mem measurableSet_Ioo ] with t ht using h_bound t ht;
        · norm_num;
      · rw [ intervalIntegrable_iff_integrableOn_Ioo_of_le ];
        · -- The function $-\frac{\log(1 - z t (1 - t))}{t}$ is integrable on $(0, 1)$ because it is continuous and bounded.
          have h_integrable : MeasureTheory.IntegrableOn (fun t => -Real.log (1 - z * t * (1 - t)) / t) (Set.Ioo 0 1) := by
            have h_cont : ContinuousOn (fun t => -Real.log (1 - z * t * (1 - t)) / t) (Set.Ioo 0 1) := by
              exact continuousOn_of_forall_continuousAt fun t ht => ContinuousAt.div ( ContinuousAt.neg ( ContinuousAt.log ( Continuous.continuousAt ( by continuity ) ) ( by nlinarith [ ht.1, ht.2, hz.1, hz.2, mul_pos hz.1 ( sub_pos.mpr ht.2 ) ] ) ) ) continuousAt_id ( ne_of_gt ht.1 )
            have h_bounded : ∀ t ∈ Set.Ioo 0 1, abs (-Real.log (1 - z * t * (1 - t)) / t) ≤ 1 / (1 - z) := by
              intros t ht
              have h_log_bound : -Real.log (1 - z * t * (1 - t)) ≤ z * t * (1 - t) / (1 - z * t * (1 - t)) := by
                have h_log_bound : ∀ x : ℝ, 0 < x ∧ x < 1 → -Real.log (1 - x) ≤ x / (1 - x) := by
                  exact fun x hx => by rw [ le_div_iff₀ ] <;> nlinarith [ Real.log_inv ( 1 - x ), Real.log_le_sub_one_of_pos ( inv_pos.mpr ( by linarith : 0 < 1 - x ) ), mul_inv_cancel₀ ( by linarith : ( 1 - x ) ≠ 0 ) ] ;
                exact h_log_bound _ ⟨ mul_pos ( mul_pos hz.1 ht.1 ) ( sub_pos.mpr ht.2 ), by nlinarith [ hz.1, hz.2, ht.1, ht.2, mul_pos hz.1 ht.1, mul_pos hz.1 ( sub_pos.mpr ht.2 ), mul_pos ht.1 ( sub_pos.mpr hz.2 ) ] ⟩;
              rw [ abs_of_nonneg ( div_nonneg ( neg_nonneg_of_nonpos ( Real.log_nonpos ( by nlinarith [ hz.1, hz.2, ht.1, ht.2, mul_pos hz.1 ht.1, mul_pos hz.1 ( sub_pos.mpr ht.2 ), mul_pos ht.1 ( sub_pos.mpr ht.2 ) ] ) ( by nlinarith [ hz.1, hz.2, ht.1, ht.2, mul_pos hz.1 ht.1, mul_pos hz.1 ( sub_pos.mpr ht.2 ), mul_pos ht.1 ( sub_pos.mpr ht.2 ) ] ) ) ) ht.1.le ) ];
              rw [ div_le_iff₀ ht.1 ];
              refine le_trans h_log_bound ?_;
              field_simp;
              rw [ div_le_div_iff₀ ] <;> nlinarith [ hz.1, hz.2, ht.1, ht.2, mul_pos hz.1 ht.1, mul_pos hz.1 ( sub_pos.mpr ht.2 ), mul_pos ht.1 ( sub_pos.mpr ht.2 ), mul_le_mul_of_nonneg_left hz.2.le ht.1.le, mul_le_mul_of_nonneg_left hz.2.le ( sub_nonneg.mpr ht.2.le ) ];
            refine' MeasureTheory.Integrable.mono' _ _ _;
            exacts [ fun t => 1 / ( 1 - z ), by norm_num, h_cont.aestronglyMeasurable measurableSet_Ioo, Filter.eventually_of_mem ( MeasureTheory.ae_restrict_mem measurableSet_Ioo ) fun t ht => h_bounded t ht ];
          convert h_integrable using 1;
        · norm_num;
    rw [ hasDerivAt_iff_tendsto_slope_zero ];
    refine' h_deriv_series.congr' _;
    filter_upwards [ self_mem_nhdsWithin, mem_nhdsWithin_of_mem_nhds ( Metric.ball_mem_nhds _ <| lt_min hz.1 <| sub_pos.2 hz.2 ) ] with h hh hh' using by rw [ h_eq h <| by simpa using hh' ] ; norm_num ; ring;

/-
The integral of 1/(1-x) from t to sqrt(t) is log(1 + sqrt(t)) for t in (0, 1).
-/
lemma inner_integral_val (t : ℝ) (ht : t ∈ Set.Ioo 0 1) : ∫ x in t..Real.sqrt t, 1 / (1 - x) = Real.log (1 + Real.sqrt t) := by
  simp +zetaDelta at *;
  rw [ integral_inv_of_pos ] <;> try nlinarith [ Real.sqrt_nonneg t, Real.sq_sqrt ht.1.le ];
  exact congrArg Real.log ( by rw [ div_eq_iff ] <;> nlinarith [ Real.sqrt_nonneg t, Real.sq_sqrt ht.1.le ] )

/-
F(x) - F(x^2) is equal to the integral of deriv F from x^2 to x for x in (0, 1).
-/
lemma F_diff_eq_integral_deriv (x : ℝ) (hx : x ∈ Set.Ioo 0 1) : F x - F (x^2) = ∫ t in x^2..x, deriv F t := by
  rw [ intervalIntegral.integral_deriv_eq_sub' ] ; aesop;
  · intro t ht;
    refine' differentiableAt_of_deriv_ne_zero _;
    -- By definition of $F$, we know that its derivative is given by the integral of $(1-t)/(1 - t^2*t*(1-t))$.
    have h_deriv : deriv F t = ∫ u in (0 : ℝ)..1, (1 - u) / (1 - t * u * (1 - u)) := by
      apply F_deriv_eq_integral;
      cases Set.mem_uIcc.mp ht <;> constructor <;> nlinarith [ hx.1, hx.2 ];
    rw [ h_deriv ];
    refine' ne_of_gt ( lt_of_lt_of_le _ ( intervalIntegral.integral_mono_on _ _ _ fun u hu => mul_le_mul_of_nonneg_left ( inv_anti₀ _ <| show 1 - t * u * ( 1 - u ) ≤ 1 by nlinarith [ hu.1, hu.2, show t * u * ( 1 - u ) ≥ 0 by exact mul_nonneg ( mul_nonneg ( show 0 ≤ t by cases Set.mem_uIcc.mp ht <;> nlinarith [ hx.1, hx.2 ] ) hu.1 ) ( sub_nonneg.mpr hu.2 ) ] ) <| sub_nonneg.mpr hu.2 ) ) <;> norm_num;
    · rw [ intervalIntegral.integral_sub ] <;> norm_num;
    · exact Continuous.intervalIntegrable ( by continuity ) _ _;
    · apply_rules [ ContinuousOn.intervalIntegrable ];
      refine' ContinuousOn.div ( continuousOn_const.sub continuousOn_id ) ( Continuous.continuousOn ( by continuity ) ) fun u hu => _;
      cases Set.mem_uIcc.mp ht <;> cases Set.mem_uIcc.mp hu <;> nlinarith [ hx.1, hx.2, mul_self_nonneg ( u - 1 / 2 ) ];
    · cases Set.mem_uIcc.mp ht <;> nlinarith [ hx.1, hx.2, hu.1, hu.2, mul_nonneg hu.1 ( sub_nonneg.mpr hu.2 ) ];
  · -- By definition of $F$, we know that its derivative is given by the integral of $(1-t)/(1-z*t*(1-t))$.
    have h_deriv : ∀ z ∈ Set.Ioo 0 1, deriv F z = ∫ t in (0 : ℝ)..1, (1 - t) / (1 - z * t * (1 - t)) := by
      exact?;
    -- Since the integrand is continuous in $z$ and the interval $[x^2, x]$ is compact, the integral is continuous on $[x^2, x]$.
    have h_cont_integrand : ContinuousOn (fun z => ∫ t in (0 : ℝ)..1, (1 - t) / (1 - z * t * (1 - t))) (Set.Icc (x^2) x) := by
      intro z hz;
      refine' intervalIntegral.tendsto_integral_filter_of_dominated_convergence _ _ _ _ _;
      use fun t => 1 / ( 1 - x );
      · filter_upwards [ self_mem_nhdsWithin ] with n hn;
        exact Measurable.aestronglyMeasurable ( by exact Measurable.mul ( measurable_const.sub measurable_id' ) ( Measurable.inv ( by exact Measurable.sub measurable_const ( by exact Measurable.mul ( measurable_const.mul measurable_id' ) ( measurable_const.sub measurable_id' ) ) ) ) );
      · refine' Filter.eventually_of_mem self_mem_nhdsWithin fun n hn => Filter.Eventually.of_forall fun t ht => _;
        rw [ Real.norm_of_nonneg ( div_nonneg ( sub_nonneg.2 <| by cases Set.mem_uIoc.1 ht <;> linarith [ hx.1, hx.2 ] ) <| sub_nonneg.2 <| by cases Set.mem_uIoc.1 ht <;> nlinarith [ hx.1, hx.2, hn.1, hn.2, mul_nonneg ( sq_nonneg x ) ( sq_nonneg t ) ] ) ] ; rw [ div_le_div_iff₀ ] <;> cases Set.mem_uIoc.1 ht <;> nlinarith [ hx.1, hx.2, hn.1, hn.2, mul_nonneg ( sq_nonneg x ) ( sq_nonneg t ) ] ;
      · norm_num;
      · refine' Filter.Eventually.of_forall fun t ht => Filter.Tendsto.div _ _ _ <;> norm_num at *;
        · exact tendsto_const_nhds;
        · exact tendsto_nhdsWithin_of_tendsto_nhds ( Continuous.tendsto' ( by continuity ) _ _ <| by norm_num );
        · nlinarith [ mul_nonneg ( sq_nonneg x ) ( sq_nonneg t ), mul_nonneg ( sq_nonneg x ) ( sq_nonneg ( 1 - t ) ), mul_nonneg ( sq_nonneg t ) ( sq_nonneg ( 1 - t ) ) ];
    rw [ Set.uIcc_of_le ( by nlinarith [ hx.1, hx.2 ] ) ] ; exact h_cont_integrand.congr fun z hz => h_deriv z ⟨ by nlinarith [ hx.1, hx.2, hz.1 ], by nlinarith [ hx.1, hx.2, hz.2 ] ⟩ ▸ rfl;

/-
The set of points (x, t) such that 0 < x < 1 and x^2 < t < x is equal to the set of points (x, t) such that 0 < t < 1 and t < x < sqrt(t).
-/
lemma region_eq : {p : ℝ × ℝ | 0 < p.1 ∧ p.1 < 1 ∧ p.1^2 < p.2 ∧ p.2 < p.1} = {p : ℝ × ℝ | 0 < p.2 ∧ p.2 < 1 ∧ p.2 < p.1 ∧ p.1 < Real.sqrt p.2} := by
  ext ⟨x, y⟩; simp [Set.mem_setOf_eq];
  constructor <;> intro h <;> refine' ⟨ _, _, _, _ ⟩ <;> nlinarith [ Real.sqrt_nonneg y, Real.sq_sqrt ( show 0 ≤ y by nlinarith ) ]

/-
S is equal to the double integral of deriv F t / (1 - x) where x goes from 0 to 1 and t goes from x^2 to x.
-/
lemma S_eq_double_integral : S_sum = ∫ x in (0 : ℝ)..1, ∫ t in x^2..x, deriv F t / (1 - x) := by
  rw [ S_eq_integral ];
  rw [ intervalIntegral.integral_of_le zero_le_one, intervalIntegral.integral_of_le zero_le_one ];
  norm_num [ MeasureTheory.integral_Ioc_eq_integral_Ioo ];
  exact MeasureTheory.setIntegral_congr_fun measurableSet_Ioo fun x hx => by rw [ F_diff_eq_integral_deriv x hx ] ;

/-
The derivative of F is bounded on the interval (0, 1).
-/
lemma deriv_F_bounded : ∃ C, ∀ z ∈ Set.Ioo 0 1, |deriv F z| ≤ C := by
  use 4 / 3;
  intro z hz
  have h_integral_bound : ∀ t ∈ Set.Ioo 0 1, |(1 - t) / (1 - z * t * (1 - t))| ≤ 4 / 3 := by
    norm_num [ abs_div ] at *;
    exact fun t ht₁ ht₂ => by rw [ div_le_iff₀ ] <;> cases abs_cases ( 1 - t ) <;> cases abs_cases ( 1 - z * t * ( 1 - t ) ) <;> nlinarith [ mul_pos hz.1 ht₁, mul_pos hz.1 ( sub_pos.2 ht₂ ), mul_pos ht₁ ( sub_pos.2 ht₂ ), pow_two_nonneg ( t - 1 / 2 ) ] ;
  have h_integral_bound : |∫ t in (0 : ℝ)..1, (1 - t) / (1 - z * t * (1 - t))| ≤ ∫ t in (0 : ℝ)..1, 4 / 3 := by
    rw [ intervalIntegral.integral_of_le zero_le_one, intervalIntegral.integral_of_le zero_le_one ];
    refine' le_trans ( MeasureTheory.norm_integral_le_integral_norm ( _ : ℝ → ℝ ) ) ( MeasureTheory.integral_mono_of_nonneg _ _ _ );
    · exact Filter.Eventually.of_forall fun x => norm_nonneg _;
    · norm_num;
    · filter_upwards [ MeasureTheory.ae_restrict_mem measurableSet_Ioc, MeasureTheory.measure_eq_zero_iff_ae_notMem.mp ( MeasureTheory.measure_singleton 1 ) ] with x hx₁ hx₂ using h_integral_bound x ⟨ hx₁.1, lt_of_le_of_ne hx₁.2 hx₂ ⟩;
  convert h_integral_bound using 1 ; norm_num [ F_deriv_eq_integral z hz ];
  norm_num

/-
The derivative of F is non-negative on (0, 1).
-/
lemma deriv_F_nonneg (z : ℝ) (hz : z ∈ Set.Ioo 0 1) : 0 ≤ deriv F z := by
  -- Apply the lemma that states the derivative of F is the integral of (1-t)/(1-zt(1-t)) over [0,1].
  have h_deriv : deriv F z = ∫ t in (0 : ℝ)..1, (1 - t) / (1 - z * t * (1 - t)) := by
    exact?;
  exact h_deriv.symm ▸ intervalIntegral.integral_nonneg ( by norm_num ) fun t ht => div_nonneg ( sub_nonneg.2 ht.2 ) ( by nlinarith [ hz.1, hz.2, ht.1, ht.2, mul_nonneg hz.1.le ( sub_nonneg.2 ht.1 ), mul_nonneg hz.1.le ( sub_nonneg.2 ht.2 ) ] ) ;

/-
The function (x, t) -> deriv F t / (1 - x) is integrable on the region 0 < x < 1, x^2 < t < x.
-/
lemma integrand_integrable_on_region :
  MeasureTheory.IntegrableOn (fun p : ℝ × ℝ => deriv F p.2 / (1 - p.1)) {p | 0 < p.1 ∧ p.1 < 1 ∧ p.1^2 < p.2 ∧ p.2 < p.1} := by
    -- Using the bounds on `deriv F`, we can show that the integrand is absolutely integrable.
    have h_integrable : MeasureTheory.IntegrableOn (fun p : ℝ × ℝ => (4 / 3) / (1 - p.1)) {p : ℝ × ℝ | 0 < p.1 ∧ p.1 < 1 ∧ p.1 ^ 2 < p.2 ∧ p.2 < p.1} := by
      have h_integrable : ∫⁻ p in {p : ℝ × ℝ | 0 < p.1 ∧ p.1 < 1 ∧ p.1^2 < p.2 ∧ p.2 < p.1}, ENNReal.ofReal (4 / 3 / (1 - p.1)) < ⊤ := by
        -- The integral of $1/(1-x)$ over $(0,1)$ is $\ln(1-x)$ evaluated from $0$ to $1$, which is $\ln(1) - \ln(0) = 0 - (-\infty) = \infty$.
        have h_integral : ∫⁻ x in Set.Ioo (0 : ℝ) 1, ENNReal.ofReal (4 / 3 / (1 - x)) * ENNReal.ofReal (x - x^2) < ⊤ := by
          -- Simplify the expression inside the integral.
          suffices h_simp : ∫⁻ x in Set.Ioo (0 : ℝ) 1, ENNReal.ofReal (4 / 3 * x) < ⊤ by
            refine' lt_of_le_of_lt ( MeasureTheory.setLIntegral_mono' measurableSet_Ioo fun x hx => _ ) h_simp;
            rw [ ← ENNReal.ofReal_mul ( div_nonneg ( by norm_num ) ( sub_nonneg.2 hx.2.le ) ) ] ; exact ENNReal.ofReal_le_ofReal ( by nlinarith [ hx.1, hx.2, mul_div_cancel₀ ( 4 / 3 ) ( by linarith [ hx.1, hx.2 ] : ( 1 - x ) ≠ 0 ) ] ) ;
          exact MeasureTheory.Integrable.lintegral_lt_top ( by exact Continuous.integrableOn_Icc ( by continuity ) |> fun h => h.mono_set <| Set.Ioo_subset_Icc_self );
        convert h_integral using 1;
        erw [ ← MeasureTheory.lintegral_indicator, ← MeasureTheory.lintegral_indicator ] <;> norm_num [ Set.indicator ];
        · erw [ MeasureTheory.lintegral_prod ];
          · congr with x ; by_cases hx : 0 < x <;> by_cases hx' : x < 1 <;> simp +decide [ hx, hx' ];
            rw [ show ( ∫⁻ y : ℝ, if x ^ 2 < y ∧ y < x then ENNReal.ofReal ( 4 / 3 / ( 1 - x ) ) else 0 ) = ∫⁻ y in Set.Ioo ( x ^ 2 ) x, ENNReal.ofReal ( 4 / 3 / ( 1 - x ) ) by rw [ ← MeasureTheory.lintegral_indicator ] <;> norm_num [ Set.indicator ] ] ; norm_num [ hx, hx' ];
          · exact Measurable.aemeasurable ( by exact Measurable.ite ( by exact MeasurableSet.inter ( measurableSet_lt measurable_const measurable_fst ) ( MeasurableSet.inter ( measurableSet_lt measurable_fst measurable_const ) ( MeasurableSet.inter ( measurableSet_lt ( measurable_fst.pow_const 2 ) measurable_snd ) ( measurableSet_lt measurable_snd measurable_fst ) ) ) ) ( by exact Measurable.ennreal_ofReal ( by exact Measurable.div ( measurable_const ) ( measurable_const.sub measurable_fst ) ) ) measurable_const );
        · exact MeasurableSet.mem ( MeasurableSet.inter ( measurableSet_lt measurable_const measurable_fst ) ( MeasurableSet.inter ( measurableSet_lt measurable_fst measurable_const ) ( MeasurableSet.inter ( measurableSet_lt ( measurable_fst.pow_const 2 ) measurable_snd ) ( measurableSet_lt measurable_snd measurable_fst ) ) ) );
      refine' ⟨ _, _ ⟩;
      · exact Measurable.aestronglyMeasurable ( by exact Measurable.div measurable_const ( measurable_const.sub measurable_fst ) );
      · rw [ MeasureTheory.hasFiniteIntegral_iff_norm ];
        refine' lt_of_le_of_lt _ h_integrable;
        refine' MeasureTheory.setLIntegral_mono' _ _;
        · exact MeasurableSet.inter ( measurableSet_lt measurable_const measurable_fst ) ( MeasurableSet.inter ( measurableSet_lt measurable_fst measurable_const ) ( MeasurableSet.inter ( measurableSet_lt ( measurable_fst.pow_const 2 ) measurable_snd ) ( measurableSet_lt measurable_snd measurable_fst ) ) );
        · exact fun p hp => by rw [ Real.norm_of_nonneg ( div_nonneg ( by norm_num ) ( by linarith [ hp.1, hp.2.1 ] ) ) ] ;
    refine' h_integrable.mono' _ _;
    · refine' Measurable.aestronglyMeasurable _;
      fun_prop;
    · rw [ MeasureTheory.ae_restrict_iff' ] <;> norm_num;
      · refine' MeasureTheory.ae_of_all _ fun p hp₁ hp₂ hp₃ hp₄ => _;
        gcongr <;> norm_num [ abs_of_pos, hp₁, hp₂, hp₃, hp₄ ];
        have h_integral_bound : ∀ t ∈ Set.Ioo 0 1, |deriv F t| ≤ 4 / 3 := by
          intro t ht; rw [ F_deriv_eq_integral t ht ] ; norm_num [ abs_div, abs_of_nonneg, ht.1.le, ht.2.le ] ; (
          rw [ intervalIntegral.integral_of_le zero_le_one ];
          rw [ abs_of_nonneg ( MeasureTheory.setIntegral_nonneg measurableSet_Ioc fun x hx => div_nonneg ( sub_nonneg.2 hx.2 ) ( by nlinarith [ hx.1, hx.2, ht.1, ht.2, mul_self_nonneg ( x - 1 / 2 ) ] ) ) ] ; refine' le_trans ( MeasureTheory.setIntegral_mono_on _ _ measurableSet_Ioc fun x hx => div_le_one_of_le₀ _ _ ) _ <;> norm_num [ ht.1.le, ht.2.le ] ; ring_nf ; (
          exact ContinuousOn.integrableOn_Icc ( by exact ContinuousOn.add ( ContinuousOn.neg ( ContinuousOn.mul continuousOn_id <| ContinuousOn.inv₀ ( Continuous.continuousOn <| by continuity ) fun x hx => by nlinarith [ hx.1, hx.2, ht.1, ht.2, mul_self_nonneg ( x - 1 / 2 ) ] ) ) <| ContinuousOn.inv₀ ( Continuous.continuousOn <| by continuity ) fun x hx => by nlinarith [ hx.1, hx.2, ht.1, ht.2, mul_self_nonneg ( x - 1 / 2 ) ] ) |> fun h => h.mono_set <| Set.Ioc_subset_Icc_self);
          · nlinarith [ hx.1, hx.2, ht.1, ht.2, mul_nonneg ht.1.le hx.1.le, mul_le_mul_of_nonneg_left ht.2.le hx.1.le ];
          · nlinarith [ ht.1, ht.2, hx.1, hx.2, mul_nonneg ht.1.le hx.1.le ]);
        exact h_integral_bound p.2 ⟨ by nlinarith, by nlinarith ⟩;
      · exact MeasurableSet.mem ( MeasurableSet.inter ( measurableSet_lt measurable_const measurable_fst ) ( MeasurableSet.inter ( measurableSet_lt measurable_fst measurable_const ) ( MeasurableSet.inter ( measurableSet_lt ( measurable_fst.pow_const 2 ) measurable_snd ) ( measurableSet_lt measurable_snd measurable_fst ) ) ) )

lemma integral_swap_equality : ∫ x in (0 : ℝ)..1, ∫ t in x^2..x, deriv F t / (1 - x) = ∫ t in (0 : ℝ)..1, ∫ x in t..Real.sqrt t, deriv F t / (1 - x) := by
  -- By Fubini's theorem, we can interchange the order of integration.
  have h_fubini : ∫ x in (0 : ℝ)..1, ∫ t in (x^2 : ℝ)..x, deriv F t / (1 - x) = ∫ t in (0 : ℝ)..1, ∫ x in (t : ℝ)..Real.sqrt t, deriv F t / (1 - x) := by
    have h_region : {p : ℝ × ℝ | 0 < p.1 ∧ p.1 < 1 ∧ p.1^2 < p.2 ∧ p.2 < p.1} = {p : ℝ × ℝ | 0 < p.2 ∧ p.2 < 1 ∧ p.2 < p.1 ∧ p.1 < Real.sqrt p.2} := by
      exact?
    have h_integrand_integrable : MeasureTheory.IntegrableOn (fun p : ℝ × ℝ => deriv F p.2 / (1 - p.1)) {p : ℝ × ℝ | 0 < p.1 ∧ p.1 < 1 ∧ p.1^2 < p.2 ∧ p.2 < p.1} := by
      exact?;
    -- By Fubini's theorem, we can interchange the order of integration since the integrand is integrable on the region.
    have h_fubini : ∫ x in Set.Ioo (0 : ℝ) 1, ∫ t in Set.Ioo (x^2 : ℝ) x, deriv F t / (1 - x) = ∫ t in Set.Ioo (0 : ℝ) 1, ∫ x in Set.Ioo (t : ℝ) (Real.sqrt t), deriv F t / (1 - x) := by
      have h_fubini : ∫ p in {p : ℝ × ℝ | 0 < p.1 ∧ p.1 < 1 ∧ p.1^2 < p.2 ∧ p.2 < p.1}, deriv F p.2 / (1 - p.1) = ∫ t in Set.Ioo (0 : ℝ) 1, ∫ x in Set.Ioo (t : ℝ) (Real.sqrt t), deriv F t / (1 - x) := by
        rw [ h_region, ← MeasureTheory.integral_indicator ];
        · erw [ MeasureTheory.integral_prod_symm ];
          · norm_num [ ← MeasureTheory.integral_indicator, Set.indicator_apply ];
            congr with x ; by_cases hx : 0 < x <;> by_cases hx' : x < 1 <;> simp +decide [ hx, hx' ];
          · rw [ ← h_region ] at *; exact MeasureTheory.integrable_indicator_iff ( by exact measurableSet_Ioi.preimage measurable_fst |> MeasurableSet.inter <| measurableSet_Iio.preimage measurable_fst |> MeasurableSet.inter <| measurableSet_lt ( measurable_fst.pow_const 2 ) measurable_snd |> MeasurableSet.inter <| measurableSet_lt measurable_snd measurable_fst ) |>.2 h_integrand_integrable;
        · exact MeasurableSet.inter ( measurableSet_lt measurable_const measurable_snd ) ( MeasurableSet.inter ( measurableSet_lt measurable_snd measurable_const ) ( MeasurableSet.inter ( measurableSet_lt measurable_snd measurable_fst ) ( measurableSet_lt measurable_fst ( Real.continuous_sqrt.measurable.comp measurable_snd ) ) ) )
      generalize_proofs at *; (
      convert h_fubini using 1
      generalize_proofs at *; (
      erw [ ← MeasureTheory.integral_indicator ] <;> norm_num [ Set.indicator ];
      erw [ ← MeasureTheory.integral_indicator ] <;> norm_num [ Set.indicator ];
      · erw [ MeasureTheory.integral_prod ];
        · congr with x ; by_cases hx : 0 < x <;> by_cases hx' : x < 1 <;> simp +decide [ hx, hx' ];
          rw [ ← MeasureTheory.integral_indicator ] <;> norm_num [ Set.indicator ];
        · rw [ ← MeasureTheory.integrable_indicator_iff ] at *;
          · convert h_integrand_integrable using 1;
            ext; simp [Set.indicator];
          · exact MeasurableSet.inter ( measurableSet_lt measurable_const measurable_fst ) ( MeasurableSet.inter ( measurableSet_lt measurable_fst measurable_const ) ( MeasurableSet.inter ( measurableSet_lt ( measurable_fst.pow_const 2 ) measurable_snd ) ( measurableSet_lt measurable_snd measurable_fst ) ) );
      · exact MeasurableSet.mem ( MeasurableSet.inter ( measurableSet_lt measurable_const measurable_fst ) ( MeasurableSet.inter ( measurableSet_lt measurable_fst measurable_const ) ( MeasurableSet.inter ( measurableSet_lt ( measurable_fst.pow_const 2 ) measurable_snd ) ( measurableSet_lt measurable_snd measurable_fst ) ) ) )));
    convert h_fubini using 1 <;> rw [ intervalIntegral.integral_of_le zero_le_one, MeasureTheory.integral_Ioc_eq_integral_Ioo ] ; norm_num [ MeasureTheory.integral_Ioc_eq_integral_Ioo ] ; ring;
    · exact MeasureTheory.setIntegral_congr_fun measurableSet_Ioo fun x hx => by rw [ ← MeasureTheory.integral_Ioc_eq_integral_Ioo, ← intervalIntegral.integral_of_le ( by nlinarith [ hx.1, hx.2 ] ) ] ; rw [ intervalIntegral.integral_mul_const ] ;
    · exact MeasureTheory.setIntegral_congr_fun measurableSet_Ioo fun x hx => by rw [ intervalIntegral.integral_of_le ( Real.le_sqrt_of_sq_le ( by nlinarith [ hx.1, hx.2 ] ) ), MeasureTheory.integral_Ioc_eq_integral_Ioo ] ;
  convert h_fubini using 1

/-
The double integral of 1/(1-x) over the region 0 < x < 1, x^2 < t < x is equal to 1/2.
-/
lemma integral_one_div_one_sub_x_finite :
  ∫ x in (0 : ℝ)..1, ∫ t in x^2..x, 1 / (1 - x) = 1 / 2 := by
    norm_num [ intervalIntegral.integral_of_le, sq ] at *;
    rw [ MeasureTheory.integral_Ioc_eq_integral_Ioo ] ; erw [ MeasureTheory.setIntegral_congr_fun measurableSet_Ioo fun x hx => by rw [ show ( x - x * x ) * ( 1 - x ) ⁻¹ = x by nlinarith [ hx.1, hx.2, mul_inv_cancel₀ ( by linarith [ hx.1, hx.2 ] : ( 1 - x ) ≠ 0 ) ] ] ] ; rw [ ← MeasureTheory.integral_Ioc_eq_integral_Ioo ] ;
    rw [ ← intervalIntegral.integral_of_le ] <;> norm_num

lemma integral_swap_equality_v2 : ∫ x in (0 : ℝ)..1, ∫ t in x^2..x, deriv F t / (1 - x) = ∫ t in (0 : ℝ)..1, ∫ x in t..Real.sqrt t, deriv F t / (1 - x) := by
  convert integral_swap_equality using 1

#check integral_swap_equality

#check S_eq_double_integral

/-
For t in (0, 1), the integral of deriv F t / (1 - x) from t to sqrt(t) is equal to deriv F t * log(1 + sqrt(t)).
-/
lemma inner_integral_eq (t : ℝ) (ht : t ∈ Set.Ioo 0 1) : ∫ x in t..Real.sqrt t, deriv F t / (1 - x) = deriv F t * Real.log (1 + Real.sqrt t) := by
  rw [show (fun x => deriv F t / (1 - x)) = (fun x => deriv F t * (1 / (1 - x))) by ext; field_simp]
  rw [intervalIntegral.integral_const_mul]
  rw [inner_integral_val t ht]

/-
For t in (0, 1), the integral of deriv F t / (1 - x) from t to sqrt(t) is equal to deriv F t * log(1 + sqrt(t)).
-/
lemma inner_integral_eq_v2 (t : ℝ) (ht : t ∈ Set.Ioo 0 1) : ∫ x in t..Real.sqrt t, deriv F t / (1 - x) = deriv F t * Real.log (1 + Real.sqrt t) := by
  rw [← inner_integral_val t ht]
  rw [← intervalIntegral.integral_const_mul]
  apply intervalIntegral.integral_congr
  intro x hx
  field_simp

/-
S is equal to the integral from 0 to 1 of deriv F t * log(1 + sqrt(t)).
-/
lemma S_eq_integral_F_deriv : S_sum = ∫ t in (0 : ℝ)..1, deriv F t * Real.log (1 + Real.sqrt t) := by
  rw [ S_eq_double_integral, integral_swap_equality ];
  rw [ intervalIntegral.integral_of_le zero_le_one ] ; rw [ MeasureTheory.integral_Ioc_eq_integral_Ioo ] ; rw [ MeasureTheory.setIntegral_congr_fun measurableSet_Ioo fun x hx => inner_integral_eq x hx ] ;
  rw [ intervalIntegral.integral_of_le zero_le_one, MeasureTheory.integral_Ioc_eq_integral_Ioo ]

/-
S is equal to the integral from 0 to 1 of deriv F (u^2) * log(1 + u) * (2 * u).
-/
lemma S_eq_integral_u : S_sum = ∫ u in (0 : ℝ)..1, deriv F (u^2) * Real.log (1 + u) * (2 * u) := by
  convert S_eq_integral_F_deriv using 1;
  symm;
  rw [ intervalIntegral.integral_of_le zero_le_one, intervalIntegral.integral_of_le zero_le_one ];
  rw [ ← MeasureTheory.integral_indicator, ← MeasureTheory.integral_indicator ] <;> norm_num [ Set.indicator ];
  -- Apply the substitution $u = \sqrt{x}$ to transform the integral.
  have h_subst : ∀ {f : ℝ → ℝ}, (∫ x in Set.Ioi (0 : ℝ), f x) = (∫ u in Set.Ioi (0 : ℝ), f (u^2) * (2 * u)) := by
    intro f; rw [ ← MeasureTheory.integral_comp_rpow_Ioi ] ; norm_num [ mul_comm ] ;
    rotate_left;
    exacts [ 2, by norm_num, by norm_num [ mul_comm ] ];
  convert @h_subst ( fun x => if 0 < x ∧ x ≤ 1 then deriv F x * Real.log ( 1 + Real.sqrt x ) else 0 ) using 1;
  · rw [ MeasureTheory.setIntegral_eq_integral_of_forall_compl_eq_zero ] ; aesop;
  · rw [ ← MeasureTheory.integral_indicator ] <;> norm_num [ Set.indicator ];
    congr with x ; split_ifs <;> simp_all +decide [ abs_of_pos, Real.sqrt_sq_eq_abs ];
    linarith

#check intervalIntegral.integral_congr_ae
#check MeasureTheory.ae_restrict_iff'

#check S_eq_integral_F_deriv

#check S_eq_integral_F_deriv

/-
S is equal to the integral from 0 to 1 of deriv F t * log(1 + sqrt(t)).
-/
lemma S_eq_integral_F_deriv_v2 : S_sum = ∫ t in (0 : ℝ)..1, deriv F t * Real.log (1 + Real.sqrt t) := by
  -- Apply the lemma S_eq_integral_F_deriv to conclude the proof.
  apply S_eq_integral_F_deriv

/-
S is equal to the integral from 0 to 1 of deriv F t * log(1 + sqrt(t)).
-/
lemma S_eq_integral_F_deriv_v3 : S_sum = ∫ t in (0 : ℝ)..1, deriv F t * Real.log (1 + Real.sqrt t) := by
  convert S_eq_integral_F_deriv using 1

/-
The derivative of G(z) is arcsin(sqrt(z)/2) / (sqrt(z) * sqrt(1 - z/4)) for z in (0, 1).
-/
noncomputable def G (z : ℝ) : ℝ := 2 * (Real.arcsin (Real.sqrt z / 2)) ^ 2

lemma G_deriv (z : ℝ) (hz : z ∈ Set.Ioo 0 1) :
  deriv G z = Real.arcsin (Real.sqrt z / 2) / (Real.sqrt z * Real.sqrt (1 - z / 4)) := by
    unfold G;
    norm_num +zetaDelta at *;
    have := @Real.hasDerivAt_arcsin ( Real.sqrt z / 2 ) ?_ ?_ <;> norm_num at *;
    · convert congr_arg ( fun x => 2 * x ) ( HasDerivAt.deriv ( this.pow 2 |> HasDerivAt.comp z <| HasDerivAt.div_const ( Real.hasDerivAt_sqrt hz.1.ne' ) _ ) ) using 1 ; norm_num ; ring;
      rw [ Real.sq_sqrt hz.1.le ];
    · linarith [ Real.sqrt_nonneg z ];
    · nlinarith [ Real.mul_self_sqrt hz.1.le ]

/-
The derivative of F is equal to the derivative of G on (0, 1).
-/
lemma deriv_F_eq_deriv_G (z : ℝ) (hz : z ∈ Set.Ioo 0 1) : deriv F z = deriv G z := by
  rw [ F_deriv_eq_integral, G_deriv z hz ];
  · -- We'll use the substitution $u = t - \frac{1}{2}$ to simplify the integral.
    have h_sub : ∫ t in (0 : ℝ)..1, (1 - t) / (1 - z * t * (1 - t)) = ∫ u in (-1 / 2 : ℝ)..1 / 2, (1 / 2 - u) / (1 - z / 4 + z * u^2) := by
      convert intervalIntegral.integral_comp_sub_right _ ( 1 / 2 ) using 2 <;> ring;
      ac_rfl;
    -- The term with $u$ in the numerator is odd, so its integral is zero.
    have h_odd : ∫ u in (-1 / 2 : ℝ)..1 / 2, (1 / 2 - u) / (1 - z / 4 + z * u^2) = ∫ u in (-1 / 2 : ℝ)..1 / 2, (1 / 2) / (1 - z / 4 + z * u^2) := by
      -- The integral of an odd function over a symmetric interval around zero is zero.
      have h_odd : ∫ u in (-1 / 2 : ℝ)..1 / 2, u / (1 - z / 4 + z * u ^ 2) = 0 := by
        have h_odd : ∀ a b : ℝ, a ≤ b → ∫ u in a..b, u / (1 - z / 4 + z * u ^ 2) = -∫ u in -b..-a, u / (1 - z / 4 + z * u ^ 2) := by
          norm_num [ ← intervalIntegral.integral_comp_neg, neg_div ];
        specialize h_odd ( -1 / 2 ) ( 1 / 2 ) ( by norm_num ) ; norm_num at * ; linarith;
      simp_all +decide [ sub_div ];
      rw [ intervalIntegral.integral_sub ( by exact Continuous.intervalIntegrable ( by exact Continuous.div continuous_const ( by continuity ) fun x => by nlinarith ) _ _ ) ( by exact Continuous.intervalIntegrable ( by exact Continuous.div ( by continuity ) ( by continuity ) fun x => by nlinarith ) _ _ ), h_odd, sub_zero ];
    -- The remaining part is $\frac{1}{2} \int_{-1/2}^{1/2} \frac{1}{z u^2 + (1-z/4)} du$.
    have h_remaining : ∫ u in (-1 / 2 : ℝ)..1 / 2, (1 / 2) / (1 - z / 4 + z * u^2) = (1 / 2) * (1 / Real.sqrt (z * (1 - z / 4))) * (Real.arctan (Real.sqrt z * (1 / 2) / Real.sqrt (1 - z / 4)) - Real.arctan (-Real.sqrt z * (1 / 2) / Real.sqrt (1 - z / 4))) := by
      rw [ intervalIntegral.integral_deriv_eq_sub' ];
      case f => exact fun u => 1 / 2 * ( 1 / Real.sqrt ( z * ( 1 - z / 4 ) ) ) * Real.arctan ( Real.sqrt z * u / Real.sqrt ( 1 - z / 4 ) );
      · grind;
      · ext u; norm_num [ mul_comm ] ; ring;
        rw [ show z + z ^ 2 * ( -1 / 4 ) = z * ( 1 + z * ( -1 / 4 ) ) by ring, Real.sqrt_mul ( by linarith [ hz.1 ] ) ] ; ring;
        norm_num [ mul_assoc, mul_comm, mul_left_comm, ne_of_gt ( Real.sqrt_pos.mpr hz.1 ), ne_of_gt ( Real.sqrt_pos.mpr ( show 0 < 1 + z * ( -1 / 4 ) by nlinarith [ hz.1, hz.2 ] ) ) ];
        rw [ Real.sq_sqrt ( by nlinarith [ hz.1, hz.2 ] ), Real.sq_sqrt ( by nlinarith [ hz.1, hz.2 ] ) ] ; ring;
        grind;
      · norm_num [ mul_comm ];
      · exact continuousOn_of_forall_continuousAt fun u hu => ContinuousAt.div continuousAt_const ( Continuous.continuousAt ( by continuity ) ) ( by nlinarith [ hz.1, hz.2 ] );
    simp_all +decide [ neg_div ];
    rw [ Real.arcsin_eq_arctan ] <;> ring <;> norm_num [ hz.1.le, hz.2.le ];
    · rw [ show z + - ( z ^ 2 * ( 1 / 4 ) ) = z * ( 1 + - ( z * ( 1 / 4 ) ) ) by ring, Real.sqrt_mul ( by linarith ) ] ; ring;
    · constructor <;> nlinarith [ Real.sqrt_nonneg z, Real.sq_sqrt hz.1.le ];
  · exact hz

/-
F is continuous at z for z in [0, 1].
-/
lemma F_continuous (z : ℝ) (hz : z ∈ Set.Icc 0 1) : ContinuousAt F z := by
  refine' ( tendsto_tsum_of_dominated_convergence _ _ _ );
  refine' fun k => 1 / ( k ^ 2 * Nat.choose ( 2 * k ) k ) * 2 ^ k;
  · -- We'll use the fact that $\binom{2k}{k} \geq 2^k$ for all $k$.
    have h_binom : ∀ k : ℕ, (Nat.choose (2 * k) k : ℝ) ≥ 2 ^ k := by
      intro k; norm_cast; induction' k with k ih <;> norm_num [ Nat.pow_succ', Nat.mul_succ, Nat.choose_succ_succ ] at *; (
      linarith [ Nat.choose_le_succ ( 2 * k ) k ]);
    -- Using the inequality $\binom{2k}{k} \geq 2^k$, we can bound the terms of the series.
    have h_bound : ∀ k : ℕ, k ≠ 0 → (1 : ℝ) / (k ^ 2 * Nat.choose (2 * k) k) * 2 ^ k ≤ 1 / k ^ 2 := by
      intro k hk; rw [ div_mul_eq_mul_div ] ; rw [ div_le_div_iff₀ ] <;> first | positivity | nlinarith [ h_binom k, show ( k : ℝ ) ^ 2 > 0 by positivity, show ( 2 : ℝ ) ^ k > 0 by positivity ] ;
    exact Summable.of_nonneg_of_le ( fun k => by positivity ) ( fun k => if hk : k = 0 then by norm_num [ hk ] else h_bound k hk ) ( Real.summable_one_div_nat_pow.2 one_lt_two );
  · exact fun k => by split_ifs <;> [ exact tendsto_const_nhds; exact Continuous.tendsto ( by continuity ) _ ] ;
  · rw [ Metric.eventually_nhds_iff ];
    refine' ⟨ 1, by norm_num, fun y hy k => _ ⟩ ; split_ifs <;> norm_num;
    · positivity;
    · -- Since $|y| \leq 2$ for $y$ close to $z$, we have $|y|^k \leq 2^k$.
      have h_abs_y_le_2 : |y| ≤ 2 := by
        exact abs_le.mpr ⟨ by linarith [ abs_lt.mp hy, hz.1, hz.2 ], by linarith [ abs_lt.mp hy, hz.1, hz.2 ] ⟩;
      exact le_trans ( mul_le_mul_of_nonneg_right ( pow_le_pow_left₀ ( abs_nonneg _ ) h_abs_y_le_2 _ ) ( by positivity ) ) ( by ring_nf; norm_num )

/-
G is continuous at z for z in [0, 1].
-/
lemma G_continuous (z : ℝ) (hz : z ∈ Set.Icc 0 1) : ContinuousAt G z := by
  refine' Continuous.continuousAt _;
  apply_rules [ Continuous.mul, Continuous.pow, Continuous.div, continuous_const, Real.continuous_arcsin.comp, Real.continuous_sqrt ]

/-
F(0) is equal to G(0).
-/
lemma F_zero_eq_G_zero : F 0 = G 0 := by
  unfold F G; norm_num;
  rw [ tsum_eq_single 0 ] <;> aesop

/-
F(z) is equal to G(z) for z in [0, 1].
-/
lemma F_eq_G (z : ℝ) (hz : z ∈ Set.Icc 0 1) : F z = G z := by
  -- Applying the Fundamental Theorem of Calculus, we can express F(z) and G(z) in terms of their derivatives.
  have h_ftc : ∀ z ∈ Set.Ioo 0 1, F z - F 0 = ∫ x in (0 : ℝ)..z, deriv F x ∧ G z - G 0 = ∫ x in (0 : ℝ)..z, deriv G x := by
    intros z hz
    constructor;
    · rw [ intervalIntegral.integral_eq_sub_of_hasDeriv_right ];
      · exact continuousOn_of_forall_continuousAt fun x hx => F_continuous x <| by constructor <;> cases Set.mem_uIcc.mp hx <;> linarith [ hz.1, hz.2 ] ;
      · norm_num +zetaDelta at *;
        intro x hx₁ hx₂; exact DifferentiableAt.hasDerivAt ( by exact differentiableAt_of_deriv_ne_zero ( by rw [ deriv_F_eq_deriv_G x ⟨ by cases hx₁ <;> linarith, by cases hx₂ <;> linarith ⟩ ] ; exact ne_of_gt ( G_deriv x ⟨ by cases hx₁ <;> linarith, by cases hx₂ <;> linarith ⟩ ▸ div_pos ( Real.arcsin_pos.mpr ( by exact div_pos ( Real.sqrt_pos.mpr ( by cases hx₁ <;> linarith ) ) zero_lt_two ) ) ( mul_pos ( Real.sqrt_pos.mpr ( by cases hx₁ <;> linarith ) ) ( Real.sqrt_pos.mpr ( by cases hx₂ <;> linarith ) ) ) ) ) ) |> HasDerivAt.hasDerivWithinAt;
      · rw [ intervalIntegrable_iff_integrableOn_Ioo_of_le hz.1.le ];
        -- Since $F$ is differentiable on $(0, 1)$, its derivative is bounded on $(0, z)$.
        have h_deriv_bounded : ∃ C, ∀ x ∈ Set.Ioo 0 z, |deriv F x| ≤ C := by
          exact ⟨ _, fun x hx => deriv_F_bounded.choose_spec x ⟨ hx.1, hx.2.trans hz.2 ⟩ ⟩;
        refine' MeasureTheory.Integrable.mono' _ _ _;
        exacts [ fun x => h_deriv_bounded.choose, Continuous.integrableOn_Icc ( by continuity ) |> fun h => h.mono_set <| Set.Ioo_subset_Icc_self, measurable_deriv F |> Measurable.aestronglyMeasurable, Filter.eventually_of_mem ( MeasureTheory.ae_restrict_mem measurableSet_Ioo ) fun x hx => h_deriv_bounded.choose_spec x hx ];
    · rw [ intervalIntegral.integral_eq_sub_of_hasDeriv_right ];
      · exact Continuous.continuousOn ( by exact Continuous.mul continuous_const <| Continuous.pow ( Real.continuous_arcsin.comp <| by continuity ) _ );
      · norm_num +zetaDelta at *;
        exact fun x hx₁ hx₂ => DifferentiableAt.hasDerivAt ( by exact differentiableAt_of_deriv_ne_zero ( by rw [ G_deriv x ⟨ by cases hx₁ <;> cases hx₂ <;> linarith, by cases hx₁ <;> cases hx₂ <;> linarith ⟩ ] ; exact div_ne_zero ( ne_of_gt ( Real.arcsin_pos.mpr ( by exact div_pos ( Real.sqrt_pos.mpr ( by cases hx₁ <;> cases hx₂ <;> linarith ) ) zero_lt_two ) ) ) ( mul_ne_zero ( ne_of_gt ( Real.sqrt_pos.mpr ( by cases hx₁ <;> cases hx₂ <;> linarith ) ) ) ( ne_of_gt ( Real.sqrt_pos.mpr ( by cases hx₁ <;> cases hx₂ <;> linarith ) ) ) ) ) ) |> HasDerivAt.hasDerivWithinAt;
      · rw [ intervalIntegrable_iff_integrableOn_Ioo_of_le hz.1.le ];
        refine' MeasureTheory.Integrable.mono' _ _ _;
        refine' fun x => Real.pi / ( Real.sqrt x * Real.sqrt ( 1 - x / 4 ) );
        · refine' MeasureTheory.Integrable.mono' _ _ _;
          refine' fun x => Real.pi / ( Real.sqrt x * Real.sqrt ( 1 - z / 4 ) );
          · -- The integral of $\frac{1}{\sqrt{x}}$ over $(0, z)$ is convergent.
            have h_sqrt_integrable : MeasureTheory.IntegrableOn (fun x => 1 / Real.sqrt x) (Set.Ioo 0 z) := by
              have h_integrable : MeasureTheory.IntegrableOn (fun x => x ^ (-1 / 2 : ℝ)) (Set.Ioo 0 z) := by
                exact ( intervalIntegral.intervalIntegrable_rpow' ( by norm_num ) ).1.mono_set ( Set.Ioo_subset_Ioc_self );
              exact h_integrable.congr_fun ( fun x hx => by norm_num [ Real.sqrt_eq_rpow, Real.rpow_neg hx.1.le ] ) measurableSet_Ioo;
            convert h_sqrt_integrable.const_mul ( Real.pi / Real.sqrt ( 1 - z / 4 ) ) using 2 ; ring;
          · exact Measurable.aestronglyMeasurable ( by exact Measurable.div measurable_const <| Measurable.mul ( Real.continuous_sqrt.measurable ) <| Real.continuous_sqrt.measurable.comp <| measurable_const.sub <| measurable_id'.div_const _ );
          · filter_upwards [ MeasureTheory.ae_restrict_mem measurableSet_Ioo ] with x hx using by rw [ Real.norm_of_nonneg ( div_nonneg Real.pi_pos.le ( mul_nonneg ( Real.sqrt_nonneg _ ) ( Real.sqrt_nonneg _ ) ) ) ] ; exact div_le_div_of_nonneg_left ( by positivity ) ( mul_pos ( Real.sqrt_pos.mpr hx.1 ) ( Real.sqrt_pos.mpr ( by linarith [ hx.2, hz.2 ] ) ) ) ( mul_le_mul_of_nonneg_left ( Real.sqrt_le_sqrt <| by linarith [ hx.2, hz.2 ] ) <| Real.sqrt_nonneg _ ) ;
        · exact ( measurable_deriv G |> Measurable.aestronglyMeasurable ) |> fun h => h.mono_measure <| MeasureTheory.Measure.restrict_le_self;
        · filter_upwards [ MeasureTheory.ae_restrict_mem measurableSet_Ioo ] with x hx ; rw [ G_deriv x ⟨ hx.1, hx.2.trans hz.2 ⟩ ] ; rw [ Real.norm_of_nonneg ( div_nonneg ( Real.arcsin_nonneg.2 <| by positivity ) <| mul_nonneg ( Real.sqrt_nonneg _ ) <| Real.sqrt_nonneg _ ) ] ; gcongr ; nlinarith [ Real.pi_pos, Real.arcsin_le_pi_div_two ( Real.sqrt x / 2 ), Real.sqrt_nonneg x, Real.sqrt_nonneg ( 1 - x / 4 ), Real.mul_self_sqrt ( show 0 ≤ x by linarith [ hx.1 ] ), Real.mul_self_sqrt ( show 0 ≤ 1 - x / 4 by linarith [ hx.2, hz.2 ] ) ] ;
  -- Since F and G are continuous on [0, 1] and differentiable on (0, 1) with equal derivatives, they must differ by a constant.
  have h_const : ∀ z ∈ Set.Ioo 0 1, F z - G z = F 0 - G 0 := by
    intros z hz
    have h_diff : ∀ x ∈ Set.Ioo 0 z, deriv F x = deriv G x := by
      exact fun x hx => deriv_F_eq_deriv_G x ⟨ hx.1, hx.2.trans hz.2 ⟩;
    simp +zetaDelta at *;
    have := h_ftc z hz.1 hz.2; rw [ intervalIntegral.integral_of_le hz.1.le, intervalIntegral.integral_of_le hz.1.le ] at this; rw [ MeasureTheory.integral_Ioc_eq_integral_Ioo, MeasureTheory.integral_Ioc_eq_integral_Ioo ] at this; rw [ MeasureTheory.setIntegral_congr_fun measurableSet_Ioo fun x hx => h_diff x hx.1 hx.2 ] at this; linarith;
  -- Since F(0) = G(0), we have F(z) - G(z) = 0 for all z in (0, 1).
  have h_zero : ∀ z ∈ Set.Ioo 0 1, F z - G z = 0 := by
    exact fun z hz => h_const z hz ▸ sub_eq_zero_of_eq ( F_zero_eq_G_zero );
  cases eq_or_lt_of_le hz.1 <;> cases eq_or_lt_of_le hz.2 <;> simp_all +decide [ sub_eq_zero ];
  · subst_vars; exact F_zero_eq_G_zero;
  · -- By continuity of $F$ and $G$ at $1$, we have $F(1) = \lim_{z \to 1^-} F(z)$ and $G(1) = \lim_{z \to 1^-} G(z)$.
    have h_cont : Filter.Tendsto F (nhdsWithin 1 (Set.Iio 1)) (nhds (F 1)) ∧ Filter.Tendsto G (nhdsWithin 1 (Set.Iio 1)) (nhds (G 1)) := by
      exact ⟨ F_continuous 1 ( by norm_num ) |> ContinuousAt.continuousWithinAt, G_continuous 1 ( by norm_num ) |> ContinuousAt.continuousWithinAt ⟩;
    exact tendsto_nhds_unique h_cont.1 ( Filter.Tendsto.congr' ( Filter.eventuallyEq_of_mem ( Ioo_mem_nhdsLT zero_lt_one ) fun x hx => by rw [ h_zero x hx.1 hx.2 ] ) h_cont.2 )

/-
The series for F_deriv_sum converges for |z| <= 1.
-/
noncomputable def F_deriv_term (k : ℕ) (z : ℝ) : ℝ := if k = 0 then 0 else z^(k-1) / ((k : ℝ) * (Nat.choose (2 * k) k : ℝ))

noncomputable def F_deriv_sum (z : ℝ) : ℝ := ∑' k : ℕ, F_deriv_term k z

lemma F_deriv_sum_converges (z : ℝ) (hz : |z| ≤ 1) : Summable (fun k => F_deriv_term k z) := by
  have h_bound : ∀ k : ℕ, k ≠ 0 → |F_deriv_term k z| ≤ 3 / 4^k := by
    intro k hk_ne_zero
    have h_binom_bound : (Nat.choose (2 * k) k : ℝ) ≥ 4 ^ k / (2 * k + 1) := by
      rw [ ge_iff_le, div_le_iff₀ ] <;> norm_cast <;> induction' k with k ih <;> norm_num [ Nat.cast_succ, Nat.mul_succ, pow_succ' ] at *;
      rcases k with ( _ | k ) <;> simp_all +decide [ Nat.choose_succ_succ, pow_succ' ];
      have := Nat.succ_mul_choose_eq ( 2 * ( k + 1 ) ) k; ( have := Nat.succ_mul_choose_eq ( 2 * ( k + 1 ) ) ( k + 1 ) ; ( norm_num [ Nat.choose_succ_succ, mul_add ] at * ; nlinarith; ) );
    unfold F_deriv_term;
    simp_all +decide [ abs_div, abs_mul ];
    rw [ div_le_div_iff₀ ] at * <;> try positivity;
    · rcases k with ( _ | k ) <;> simp_all +decide [ pow_succ' ];
      rw [ div_le_iff₀ ] at h_binom_bound <;> nlinarith [ pow_le_pow_of_le_one ( abs_nonneg z ) hz ( show k ≥ 0 by positivity ), pow_pos ( by positivity : ( 0 : ℝ ) < 4 ) k ];
    · exact mul_pos ( Nat.cast_pos.mpr ( Nat.pos_of_ne_zero hk_ne_zero ) ) ( Nat.cast_pos.mpr ( Nat.choose_pos ( by linarith ) ) );
  have h_summable : Summable (fun k : ℕ => 3 / (4 : ℝ)^k) := by
    simpa using summable_geometric_of_lt_one ( by norm_num ) ( by norm_num : ( 4 : ℝ ) ⁻¹ < 1 ) |> Summable.mul_left 3;
  -- Apply the comparison test using the bound we have and the summability of the series ∑' k : ℕ, 3 / (4 : ℝ)^k.
  have h_comparison : Summable (fun k : ℕ => |F_deriv_term k z|) := by
    rw [ ← summable_nat_add_iff 1 ] at *;
    exact Summable.of_nonneg_of_le ( fun n => abs_nonneg _ ) ( fun n => h_bound _ ( Nat.succ_ne_zero _ ) ) h_summable;
  exact h_comparison.of_abs

/-
F_deriv_sum is continuous on [0, 1].
-/
lemma F_deriv_sum_continuous : ContinuousOn F_deriv_sum (Set.Icc 0 1) := by
  have h_uniform : Summable (fun k : ℕ => (1 : ℝ) / (k * (Nat.choose (2 * k) k))) := by
    -- We can compare our series with the convergent p-series $\sum_{k=1}^{\infty} \frac{1}{k^2}$.
    have h_comparison : ∀ k : ℕ, k ≥ 1 → (1 : ℝ) / (k * Nat.choose (2 * k) k) ≤ 1 / (k ^ 2) := by
      intro k hk; gcongr ; norm_cast ; induction hk <;> norm_num [ Nat.pow_succ', Nat.mul_succ, Nat.choose_succ_succ ] at *;
      nlinarith [ Nat.choose_pos ( by linarith : ‹_› ≤ 2 * ‹_› + 1 ) ];
    exact Summable.of_nonneg_of_le ( fun k => by positivity ) ( fun k => if hk : k ≥ 1 then h_comparison k hk else by interval_cases k ; norm_num ) ( Real.summable_one_div_nat_pow.2 one_lt_two );
  refine' continuousOn_tsum _ _ _;
  use fun k => 1 / ( k * Nat.choose ( 2 * k ) k );
  · intro k; exact Continuous.continuousOn ( by unfold F_deriv_term; split_ifs <;> continuity ) ;
  · exact?;
  · unfold F_deriv_term;
    intro n x hx; split_ifs <;> norm_num;
    · positivity;
    · exact le_of_le_of_eq ( mul_le_mul_of_nonneg_right ( pow_le_one₀ ( abs_nonneg x ) ( abs_le.mpr ⟨ by linarith [ hx.1 ], by linarith [ hx.2 ] ⟩ ) ) ( by positivity ) ) ( by ring )

/-
The series for F_deriv_sum converges for |z| < 4.
-/
lemma F_deriv_sum_converges_large (z : ℝ) (hz : |z| < 4) : Summable (fun k => F_deriv_term k z) := by
  have h_comparison : ∀ k : ℕ, 1 ≤ k → |F_deriv_term k z| ≤ |z|^(k-1) * (4 / 4^k) := by
    intro k hk
    have h_binom : (Nat.choose (2 * k) k : ℝ) ≥ 4^k / (2 * k + 1) := by
      rw [ ge_iff_le, div_le_iff₀ ] <;> norm_cast <;> induction hk <;> norm_num [ Nat.mul_succ, pow_succ' ] at *;
      have := Nat.succ_mul_choose_eq ( 2 * ‹_› ) ‹_›; ( have := Nat.succ_mul_choose_eq ( 2 * ‹_› + 1 ) ‹_›; ( norm_num [ Nat.choose_succ_succ ] at * ; nlinarith; ) );
    unfold F_deriv_term
    simp [h_binom];
    split_ifs <;> simp_all +decide [ abs_div, abs_mul, abs_of_nonneg, add_nonneg ];
    rw [ div_le_iff₀ ] at * <;> try positivity;
    · rw [ mul_assoc ];
      exact le_mul_of_one_le_right ( by positivity ) ( by rw [ div_mul_eq_mul_div, le_div_iff₀ ] <;> first | positivity | nlinarith [ show ( k : ℝ ) ≥ 1 by norm_cast ] );
    · exact mul_pos ( by positivity ) ( Nat.cast_pos.mpr ( Nat.choose_pos ( by linarith ) ) );
  have h_series_converges : Summable (fun k : ℕ => |z|^(k - 1) * (4 / 4^k)) := by
    rw [ ← summable_nat_add_iff 1 ] ; ring_nf ; norm_num [ pow_mul', hz ];
    simpa only [ ← mul_pow ] using summable_geometric_of_lt_one ( by positivity ) ( by linarith );
  rw [ ← summable_nat_add_iff 1 ] at *;
  exact Summable.of_norm <| h_series_converges.of_nonneg_of_le ( fun n => norm_nonneg _ ) fun n => h_comparison _ le_add_self

/-
F has derivative F_deriv_sum at z for |z| < 3.
-/
noncomputable def F_term (k : ℕ) (z : ℝ) : ℝ := if k = 0 then 0 else z^k / ((k : ℝ)^2 * (Nat.choose (2 * k) k : ℝ))

lemma F_hasDerivAt (z : ℝ) (hz : |z| < 3) : HasDerivAt F (F_deriv_sum z) z := by
  have := @F_deriv_sum_converges_large;
  convert hasDerivAt_tsum_of_isPreconnected _ _ _ using 1;
  rotate_left;
  exact ℕ;
  exact ℝ;
  exact ℝ;
  all_goals try infer_instance;
  use fun n => 3 ^ ( n - 1 ) / ( n * Nat.choose ( 2 * n ) n );
  use fun n z => if n = 0 then 0 else z ^ n / ( n ^ 2 * Nat.choose ( 2 * n ) n );
  use fun n z => if n = 0 then 0 else z ^ ( n - 1 ) / ( n * Nat.choose ( 2 * n ) n );
  exact Set.Ioo ( -3 ) 3;
  exact 0;
  exact z;
  · convert this 3 ( by norm_num ) using 1;
    ext ( _ | n ) <;> norm_num [ F_deriv_term ];
  · exact isOpen_Ioo;
  · exact isPreconnected_Ioo;
  · constructor <;> intro h;
    · aesop;
    · convert h _ _ _ _ _ using 1;
      · intro n y hy; by_cases hn : n = 0 <;> simp +decide [ hn ];
        · exact hasDerivAt_const _ _;
        · convert HasDerivAt.div_const ( hasDerivAt_pow n y ) _ using 1 ; ring;
          simp +decide [ sq, mul_assoc, hn ];
      · intro n y hy; split_ifs <;> norm_num;
        · positivity;
        · gcongr ; exact abs_le.mpr ⟨ by linarith [ hy.1 ], by linarith [ hy.2 ] ⟩;
      · norm_num;
      · exact ⟨ _, hasSum_single 0 fun n hn => by aesop ⟩;
      · exact abs_lt.mp hz

/-
The integral of deriv F(u^2) * log(1+u) * 2u from 0 to 1 is equal to the integral of deriv F(4sin^2(theta)) * log(1+2sin(theta)) * 8sin(theta)cos(theta) from 0 to pi/6.
-/
lemma S_integral_substitution :
  ∫ u in (0 : ℝ)..1, deriv F (u^2) * Real.log (1 + u) * (2 * u) =
  ∫ theta in (0 : ℝ)..Real.pi / 6, deriv F (4 * Real.sin theta ^ 2) * Real.log (1 + 2 * Real.sin theta) * 8 * Real.sin theta * Real.cos theta := by
    have int_comparison : ∫ u in (0 : ℝ)..1, deriv F (u ^ 2) * Real.log (1 + u) * (2 * u) = ∫ u in (0 : ℝ)..Real.pi / 6, deriv F ((2 * Real.sin u) ^ 2) * Real.log (1 + 2 * Real.sin u) * (2 * (2 * Real.sin u)) * 2 * Real.cos u := by
      have int_comparison : ∫ u in (0 : ℝ)..1, deriv F (u ^ 2) * Real.log (1 + u) * (2 * u) = ∫ u in (Set.Icc 0 (Real.pi / 6)), deriv F ((2 * Real.sin u) ^ 2) * Real.log (1 + 2 * Real.sin u) * (2 * (2 * Real.sin u)) * 2 * Real.cos u := by
        have int_comparison : ∫ u in (0 : ℝ)..1, deriv F (u ^ 2) * Real.log (1 + u) * (2 * u) = ∫ u in (Set.image (fun u => 2 * Real.sin u) (Set.Icc 0 (Real.pi / 6))), deriv F (u ^ 2) * Real.log (1 + u) * (2 * u) := by
          rw [ show ( fun u => 2 * Real.sin u ) '' Set.Icc 0 ( Real.pi / 6 ) = Set.Icc 0 1 from ?_ ];
          · rw [ MeasureTheory.integral_Icc_eq_integral_Ioc, intervalIntegral.integral_of_le zero_le_one ];
          · ext x
            simp [Set.mem_image];
            constructor;
            · rintro ⟨ y, ⟨ hy₁, hy₂ ⟩, rfl ⟩ ; exact ⟨ by linarith [ Real.sin_nonneg_of_nonneg_of_le_pi hy₁ ( by linarith ) ], by linarith [ show Real.sin y ≤ 1 / 2 by rw [ ← Real.cos_pi_div_two_sub ] ; rw [ ← Real.cos_pi_div_three ] ; exact Real.cos_le_cos_of_nonneg_of_le_pi ( by linarith ) ( by linarith ) ( by linarith ) ] ⟩ ;
            · intro hx
              use Real.arcsin (x / 2);
              rw [ Real.arcsin_le_iff_le_sin ] <;> norm_num;
              · exact ⟨ ⟨ by linarith, by linarith ⟩, by rw [ Real.sin_arcsin ] <;> linarith ⟩;
              · constructor <;> linarith;
              · constructor <;> linarith [ Real.pi_pos ];
        rw [ int_comparison, MeasureTheory.integral_image_eq_integral_abs_deriv_smul ] <;> norm_num [ mul_assoc ];
        any_goals intro x hx₁ hx₂; exact HasDerivAt.hasDerivWithinAt ( by simpa using HasDerivAt.const_mul 2 ( Real.hasDerivAt_sin x ) );
        · exact MeasureTheory.setIntegral_congr_fun measurableSet_Icc fun x hx => by rw [ abs_of_nonneg ( mul_nonneg zero_le_two ( Real.cos_nonneg_of_mem_Icc ⟨ by linarith [ Real.pi_pos, hx.1 ], by linarith [ Real.pi_pos, hx.2 ] ⟩ ) ) ] ; ring;
        · exact fun x hx y hy hxy => Real.injOn_sin ⟨ by linarith [ Real.pi_pos, hx.1 ], by linarith [ Real.pi_pos, hx.2 ] ⟩ ⟨ by linarith [ Real.pi_pos, hy.1 ], by linarith [ Real.pi_pos, hy.2 ] ⟩ <| by linarith;
      rw [ int_comparison, MeasureTheory.integral_Icc_eq_integral_Ioc, intervalIntegral.integral_of_le ( by positivity ) ];
    exact int_comparison.trans ( by congr; ext; ring )

/-
The integrand simplifies to 4 * theta * log(1 + 2sin(theta)) for theta in (0, pi/6).
-/
noncomputable def S_simplified_integrand (theta : ℝ) : ℝ := 4 * theta * Real.log (1 + 2 * Real.sin theta)

lemma integrand_simplification (theta : ℝ) (htheta : theta ∈ Set.Ioo 0 (Real.pi / 6)) :
  deriv F (4 * Real.sin theta ^ 2) * Real.log (1 + 2 * Real.sin theta) * 8 * Real.sin theta * Real.cos theta = S_simplified_integrand theta := by
    unfold S_simplified_integrand;
    rw [ show deriv F ( 4 * Real.sin theta ^ 2 ) = deriv G ( 4 * Real.sin theta ^ 2 ) from deriv_F_eq_deriv_G _ ?_ ];
    · rw [ G_deriv ] <;> norm_num;
      · rw [ Real.sqrt_sq ( le_of_lt ( Real.sin_pos_of_pos_of_lt_pi htheta.1 ( by linarith [ Real.pi_pos, htheta.2 ] ) ) ) ] ; rw [ ← Real.cos_sq' ] ; rw [ Real.sqrt_sq ( le_of_lt ( Real.cos_pos_of_mem_Ioo ⟨ by linarith [ Real.pi_pos, htheta.1 ], by linarith [ Real.pi_pos, htheta.2 ] ⟩ ) ) ] ; ring;
        rw [ Real.arcsin_sin ] <;> try linarith [ htheta.1, htheta.2 ];
        simp +decide [ mul_assoc, mul_comm, mul_left_comm, ne_of_gt ( Real.sin_pos_of_pos_of_lt_pi htheta.1 ( by linarith [ Real.pi_pos, htheta.2 ] ) ), ne_of_gt ( Real.cos_pos_of_mem_Ioo ⟨ by linarith [ Real.pi_pos, htheta.1 ], by linarith [ Real.pi_pos, htheta.2 ] ⟩ : 0 < Real.cos theta ) ];
      · exact ⟨ pow_pos ( Real.sin_pos_of_pos_of_lt_pi htheta.1 ( by linarith [ Real.pi_pos, htheta.2 ] ) ) 2, by nlinarith [ show Real.sin theta < 1 / 2 from by rw [ ← Real.cos_pi_div_two_sub ] ; rw [ ← Real.cos_pi_div_three ] ; exact Real.cos_lt_cos_of_nonneg_of_le_pi ( by linarith [ Real.pi_pos, htheta.1, htheta.2 ] ) ( by linarith [ Real.pi_pos, htheta.1, htheta.2 ] ) ( by linarith [ Real.pi_pos, htheta.1, htheta.2 ] ), Real.sin_pos_of_pos_of_lt_pi htheta.1 ( by linarith [ Real.pi_pos, htheta.2 ] ) ] ⟩;
    · simp +zetaDelta at *;
      exact ⟨ pow_pos ( Real.sin_pos_of_pos_of_lt_pi htheta.1 ( by linarith ) ) _, by nlinarith [ show Real.sin theta < 1 / 2 by rw [ ← Real.cos_pi_div_two_sub ] ; rw [ ← Real.cos_pi_div_three ] ; exact Real.cos_lt_cos_of_nonneg_of_le_pi ( by linarith ) ( by linarith ) ( by linarith ), Real.sin_pos_of_pos_of_lt_pi htheta.1 ( by linarith ) ] ⟩

/-
S is equal to S_integral_theta.
-/
noncomputable def S_integral_theta : ℝ := ∫ theta in (0 : ℝ)..Real.pi / 6, deriv F (4 * Real.sin theta ^ 2) * Real.log (1 + 2 * Real.sin theta) * 8 * Real.sin theta * Real.cos theta

lemma S_eq_S_integral_theta : S_sum = S_integral_theta := by
  convert S_eq_integral_u using 1 ; unfold S_integral_theta ; ring;
  convert S_integral_substitution.symm using 2 ; ring;
  · ring;
  · exact funext fun x => by ring;

/-
1 + 2sin(theta) is equal to 4sin(theta/2 + pi/12)sin(theta/2 + 5pi/12).
-/
lemma sin_factorization (theta : ℝ) :
  1 + 2 * Real.sin theta = 4 * Real.sin (theta / 2 + Real.pi / 12) * Real.sin (theta / 2 + 5 * Real.pi / 12) := by
    rw [ show 5 * Real.pi / 12 = Real.pi / 4 + Real.pi / 6 by ring, show Real.pi / 12 = Real.pi / 4 - Real.pi / 6 by ring ] ; norm_num [ Real.sin_add, Real.sin_sub ] ; ring;
    rw [ show Real.pi * ( 5 / 12 ) = Real.pi / 4 + Real.pi / 6 by ring, show Real.pi * ( 1 / 12 ) = Real.pi / 4 - Real.pi / 6 by ring ] ; norm_num [ Real.sin_add, Real.sin_sub, Real.cos_add, Real.cos_sub ] ; ring;
    rw [ show theta = 2 * ( theta / 2 ) by ring ] ; norm_num [ Real.sin_two_mul, Real.cos_two_mul ] ; ring; norm_num [ Real.sin_sq, Real.cos_sq ] ; ring;

/-
The real part of log(1 - exp(2*I*phi)) is log(2*sin(phi)) for phi in (0, pi).
-/
lemma real_part_neg_log_one_sub_exp_mul_two_I_phi (phi : ℝ) (hphi : phi ∈ Set.Ioo 0 Real.pi) :
  (Complex.log (1 - Complex.exp (2 * Complex.I * phi))).re = Real.log (2 * Real.sin phi) := by
    norm_num [ Complex.log_re, Complex.log_im, Complex.normSq, Complex.norm_def, Complex.exp_re, Complex.exp_im ] at *;
    ring_nf; norm_num [ Real.sin_sq, Real.cos_sq ] ; ring; (
    rw [ show 2 - Real.cos ( phi * 2 ) * 2 = ( Real.sin phi * 2 ) ^ 2 by rw [ mul_pow, Real.sin_sq, Real.cos_sq ] ; ring ] ; rw [ Real.sqrt_sq ( mul_nonneg ( Real.sin_nonneg_of_nonneg_of_le_pi hphi.1.le hphi.2.le ) zero_le_two ) ] ;);

/-
The partial sums of z^n are bounded for |z|=1, z != 1.
-/
lemma partial_sums_z_pow_bounded (z : ℂ) (hz : ‖z‖ = 1) (hz1 : z ≠ 1) :
  ∃ C, ∀ N, ‖∑ n ∈ Finset.range N, z ^ n‖ ≤ C := by
    -- The partial sums are $\frac{1-z^N}{1-z}$.
    have h_partial_sums : ∀ N : ℕ, ∑ n ∈ Finset.range N, z^n = (1 - z^N) / (1 - z) := by
      exact fun N => by rw [ ← neg_div_neg_eq, geom_sum_eq ] <;> aesop;
    norm_num [ h_partial_sums ];
    exact ⟨ 2 / ‖1 - z‖, fun N => div_le_div_of_nonneg_right ( le_trans ( norm_sub_le _ _ ) ( by norm_num [ hz ] ) ) ( norm_nonneg _ ) ⟩

/-
Dirichlet's test for convergence of a series.
-/
lemma dirichlet_test_convergence {a : ℕ → ℂ} {b : ℕ → ℝ}
  (h_bounded : ∃ C, ∀ N, ‖∑ n ∈ Finset.range N, a n‖ ≤ C)
  (h_mono : Antitone b)
  (h_lim : Filter.Tendsto b Filter.atTop (nhds 0)) :
  ∃ L, Filter.Tendsto (fun N => ∑ n ∈ Finset.range N, a n * b n) Filter.atTop (nhds L) := by
    -- Using summation by parts, we can rewrite the partial sums.
    have h_summation_parts : ∀ N : ℕ, ∑ n ∈ Finset.range N, a n * b n = (∑ n ∈ Finset.range N, a n) * b (N - 1) + ∑ n ∈ Finset.range (N - 1), (∑ k ∈ Finset.range (n + 1), a k) * (b n - b (n + 1)) := by
      intro N; induction N <;> simp_all +decide [ Finset.sum_range_succ ] ; ring;
      cases ‹ℕ› <;> norm_num [ Finset.sum_range_succ ] ; ring;
    -- The series $\sum_{n=0}^{\infty} (\sum_{k=0}^{n} a_k) (b_n - b_{n+1})$ converges absolutely.
    have h_abs_conv : Summable (fun n => ‖(∑ k ∈ Finset.range (n + 1), a k) * (b n - b (n + 1))‖) := by
      -- Since $|S_n|$ is bounded and $|b_n - b_{n+1}|$ is non-negative and decreasing, their product is absolutely summable.
      have h_abs_summable : Summable (fun n => ‖(∑ k ∈ Finset.range (n + 1), a k)‖ * (b n - b (n + 1))) := by
        have h_abs_summable : Summable (fun n => (b n - b (n + 1))) := by
          rw [ summable_iff_not_tendsto_nat_atTop_of_nonneg ];
          · -- The series $\sum_{n=0}^{N-1} (b_n - b_{n+1})$ is a telescoping series.
            have h_telescoping : ∀ N : ℕ, ∑ n ∈ Finset.range N, (b n - b (n + 1)) = b 0 - b N := by
              exact fun N => by rw [ Finset.sum_range_sub' ] ;
            exact fun h => not_tendsto_atTop_of_tendsto_nhds ( by simpa only [ h_telescoping ] using tendsto_const_nhds.sub h_lim ) h;
          · exact fun n => sub_nonneg_of_le <| h_mono <| Nat.le_succ _;
        exact Summable.of_nonneg_of_le ( fun n => mul_nonneg ( norm_nonneg _ ) ( sub_nonneg.mpr ( h_mono ( Nat.le_succ _ ) ) ) ) ( fun n => mul_le_mul_of_nonneg_right ( h_bounded.choose_spec _ ) ( sub_nonneg.mpr ( h_mono ( Nat.le_succ _ ) ) ) ) ( h_abs_summable.mul_left _ );
      convert h_abs_summable using 2 ; norm_cast ; norm_num [ abs_mul, abs_of_nonneg, h_mono ( Nat.le_succ _ ) ];
      exact Or.inl <| mod_cast abs_of_nonneg <| sub_nonneg_of_le <| h_mono <| Nat.le_succ _;
    -- The series $\sum_{n=0}^{\infty} (\sum_{k=0}^{n} a_k) (b_n - b_{n+1})$ converges.
    have h_conv : Filter.Tendsto (fun N => ∑ n ∈ Finset.range N, (∑ k ∈ Finset.range (n + 1), a k) * (b n - b (n + 1))) Filter.atTop (nhds (∑' n, (∑ k ∈ Finset.range (n + 1), a k) * (b n - b (n + 1)))) := by
      exact h_abs_conv.of_norm.hasSum.tendsto_sum_nat;
    -- The term $(\sum_{n=0}^{N-1} a_n) b_{N-1}$ tends to $0$ as $N$ tends to infinity.
    have h_term_zero : Filter.Tendsto (fun N => (∑ n ∈ Finset.range N, a n) * b (N - 1)) Filter.atTop (nhds 0) := by
      exact squeeze_zero_norm ( fun N => by simpa [ abs_mul ] using mul_le_mul_of_nonneg_right ( h_bounded.choose_spec N ) ( by positivity ) ) ( by simpa using Filter.Tendsto.mul ( tendsto_const_nhds ) ( h_lim.norm.comp ( Filter.tendsto_sub_atTop_nat 1 ) ) );
    exact ⟨ _, by simpa only [ h_summation_parts, zero_add ] using h_term_zero.add ( h_conv.comp ( Filter.tendsto_sub_atTop_nat 1 ) ) ⟩

/-
The partial sums of z^n/n converge for |z|=1, z != 1.
-/
lemma tendsto_sum_z_pow_n_div_n_convergence (z : ℂ) (hz : ‖z‖ = 1) (hz1 : z ≠ 1) :
  ∃ L, Filter.Tendsto (fun N => ∑ n ∈ Finset.range N, if n = 0 then 0 else z ^ n / n) Filter.atTop (nhds L) := by
    -- Apply Dirichlet's test with the partial sums of $z^n$ and the sequence $1/n$.
    have h_dirichlet : ∃ C, ∀ N, ‖∑ n ∈ Finset.range N, z ^ n‖ ≤ C ∧ Antitone (fun n : ℕ => (1 : ℝ) / (n + 1)) ∧ Filter.Tendsto (fun n : ℕ => (1 : ℝ) / (n + 1)) Filter.atTop (nhds 0) := by
      obtain ⟨ C, hC ⟩ := partial_sums_z_pow_bounded z hz hz1;
      exact ⟨ C, fun N => ⟨ hC N, fun n m hnm => by simpa using inv_anti₀ ( by positivity ) ( by norm_cast; linarith ), tendsto_one_div_add_atTop_nhds_zero_nat ⟩ ⟩;
    -- Apply Dirichlet's test with the partial sums of $z^n$ and the sequence $1/n$ to conclude the existence of the limit.
    obtain ⟨C, hC⟩ := h_dirichlet;
    have h_apply_dirichlet : ∃ L : ℂ, Filter.Tendsto (fun N => ∑ n ∈ Finset.range N, z ^ (n + 1) / (n + 1 : ℝ)) Filter.atTop (nhds L) := by
      have h_apply_dirichlet : ∃ L : ℂ, Filter.Tendsto (fun N => ∑ n ∈ Finset.range N, (z ^ (n + 1)) * (1 / (n + 1 : ℝ))) Filter.atTop (nhds L) := by
        have h_bounded : ∃ C, ∀ N, ‖∑ n ∈ Finset.range N, z ^ (n + 1)‖ ≤ C := by
          simp_all +decide [ pow_succ', ← Finset.mul_sum _ _ _, ← Finset.sum_mul ];
          exact ⟨ C, fun N => hC N |>.1 ⟩
        convert dirichlet_test_convergence _ _ _;
        rotate_left;
        exacts [ fun n => 1 / ( n + 1 : ℝ ), h_bounded, hC 0 |>.2.1, hC 0 |>.2.2, by norm_num ];
      aesop;
    obtain ⟨ L, hL ⟩ := h_apply_dirichlet; use L; rw [ ← Filter.tendsto_add_atTop_iff_nat 1 ] ; simp_all +decide [ Finset.sum_range_succ' ] ;

/-
The partial sums of z^n/n converge to -log(1-z) for |z|=1, z != 1.
-/
lemma tendsto_sum_z_pow_n_div_n (z : ℂ) (hz : ‖z‖ = 1) (hz1 : z ≠ 1) :
  Filter.Tendsto (fun N => ∑ n ∈ Finset.range N, if n = 0 then 0 else z ^ n / n) Filter.atTop (nhds (-Complex.log (1 - z))) := by
    have := @Complex.tendsto_tsum_powerSeries_nhdsWithin_lt;
    -- We'll use the fact that if the series $\sum_{n=1}^{\infty} \frac{z^n}{n}$ converges, then its sum is $-\log(1-z)$.
    have h_sum : ∀ r : ℝ, 0 < r ∧ r < 1 → ∑' n : ℕ, (if n = 0 then 0 else z^n / (n : ℂ)) * r^n = -Complex.log (1 - r * z) := by
      intro r hr
      have h_series : ∀ r : ℝ, 0 < r ∧ r < 1 → HasSum (fun n : ℕ => (r * z) ^ n / (n : ℂ)) (-Complex.log (1 - r * z)) := by
        intro r hr
        have h_series : ∀ r : ℝ, 0 < r ∧ r < 1 → HasSum (fun n : ℕ => (r * z) ^ n / (n : ℂ)) (-Complex.log (1 - r * z)) := by
          intro r hr
          have h_abs : ‖r * z‖ < 1 := by
            simpa [ abs_of_pos hr.1, hz ] using hr.2
          exact?;
        exact h_series r hr;
      convert HasSum.tsum_eq ( h_series r hr ) using 1;
      exact tsum_congr fun n => by by_cases hn : n = 0 <;> simp +decide [ hn, div_eq_mul_inv, mul_pow, mul_assoc, mul_comm, mul_left_comm ] ;
    -- By the properties of the logarithm and the fact that $|z| = 1$, we have $\lim_{r \to 1^-} -\log(1 - rz) = -\log(1 - z)$.
    have h_log_limit : Filter.Tendsto (fun r : ℝ => -Complex.log (1 - r * z)) (nhdsWithin 1 (Set.Iio 1)) (nhds (-Complex.log (1 - z))) := by
      refine' Filter.Tendsto.neg _;
      refine' Filter.Tendsto.comp ( Complex.differentiableAt_log _ |> DifferentiableAt.continuousAt ) _;
      · norm_num [ Complex.slitPlane ];
        contrapose! hz1; simp_all +decide [ Complex.ext_iff ] ;
        nlinarith [ Complex.normSq_apply z, Complex.sq_norm z ];
      · exact tendsto_nhdsWithin_of_tendsto_nhds ( Continuous.tendsto' ( by continuity ) _ _ <| by norm_num );
    contrapose! this;
    use fun n => if n = 0 then 0 else z^n / (n : ℂ);
    obtain ⟨ l, hl ⟩ := tendsto_sum_z_pow_n_div_n_convergence z hz hz1;
    refine' ⟨ l, hl, _ ⟩;
    intro H;
    have := H.congr' ( Filter.eventuallyEq_of_mem ( Ioo_mem_nhdsLT zero_lt_one ) fun x hx => h_sum x hx );
    have := tendsto_nhds_unique this h_log_limit; aesop;

/-
The partial sums of -cos(2n*phi)/n converge to log(2sin(phi)) for phi in (0, pi).
-/
lemma log_two_sin_fourier_series_limit (phi : ℝ) (hphi : phi ∈ Set.Ioo 0 Real.pi) :
  Filter.Tendsto (fun N => - ∑ n ∈ Finset.range N, if n = 0 then 0 else Real.cos (2 * n * phi) / n) Filter.atTop (nhds (Real.log (2 * Real.sin phi))) := by
    have h_real_part : Filter.Tendsto (fun N => ∑ n ∈ Finset.range N, if n = 0 then 0 else (Complex.exp (2 * Complex.I * phi)) ^ n / (n : ℂ)) Filter.atTop (nhds (-Complex.log (1 - Complex.exp (2 * Complex.I * phi)))) := by
      convert tendsto_sum_z_pow_n_div_n ( Complex.exp ( 2 * Complex.I * phi ) ) ?_ ?_ using 1 <;> norm_num [ Complex.norm_exp ];
      rw [ Complex.exp_eq_one_iff ];
      norm_num [ Complex.ext_iff, hphi.1.ne', hphi.2.ne ];
      rintro ⟨ _ | k ⟩ <;> norm_num <;> nlinarith [ hphi.1, hphi.2 ];
    convert Complex.continuous_re.continuousAt.tendsto.comp ( h_real_part.neg ) using 2 ; norm_num [ Complex.norm_exp ];
    · congr! 1;
      split_ifs <;> simp +decide [ *, ← Complex.exp_nat_mul, Complex.exp_re ] ; ring;
    · convert real_part_neg_log_one_sub_exp_mul_two_I_phi phi hphi |> Eq.symm using 1 ; ring

/-
log(1 + 2sin(theta)) decomposes into two logs.
-/
lemma log_one_plus_two_sin_decomposition (theta : ℝ) (htheta : theta ∈ Set.Icc 0 (Real.pi / 6)) :
  Real.log (1 + 2 * Real.sin theta) =
  Real.log (2 * Real.sin (theta / 2 + Real.pi / 12)) + Real.log (2 * Real.sin (theta / 2 + 5 * Real.pi / 12)) := by
    rw [ ← Real.log_mul ];
    · convert congr_arg _ ( sin_factorization theta ) using 2 ; ring;
    · exact ne_of_gt ( mul_pos zero_lt_two ( Real.sin_pos_of_pos_of_lt_pi ( by linarith [ Real.pi_pos, htheta.1 ] ) ( by linarith [ Real.pi_pos, htheta.2 ] ) ) );
    · exact mul_ne_zero two_ne_zero ( ne_of_gt ( Real.sin_pos_of_pos_of_lt_pi ( by linarith [ Real.pi_pos, htheta.1 ] ) ( by linarith [ Real.pi_pos, htheta.2 ] ) ) )

/-
cos(n(theta + pi/6)) + cos(n(theta + 5pi/6)) = 2cos(n*pi/3)cos(n(theta + pi/2)).
-/
lemma cos_sum_identity (n : ℕ) (theta : ℝ) :
  Real.cos (n * (theta + Real.pi / 6)) + Real.cos (n * (theta + 5 * Real.pi / 6)) =
  2 * Real.cos (n * Real.pi / 3) * Real.cos (n * (theta + Real.pi / 2)) := by
    rw [ Real.cos_add_cos ] ; ring;
    norm_num

/-
The partial sums of the cosine series converge to log(1 + 2sin(theta)).
-/
lemma log_one_plus_two_sin_series_limit (theta : ℝ) (htheta : theta ∈ Set.Icc 0 (Real.pi / 6)) :
  Filter.Tendsto (fun N => - ∑ n ∈ Finset.range N, if n = 0 then 0 else (Real.cos (n * (theta + Real.pi / 6)) + Real.cos (n * (theta + 5 * Real.pi / 6))) / n) Filter.atTop (nhds (Real.log (1 + 2 * Real.sin theta))) := by
    -- Apply the decomposition of log(1 + 2sin(theta)) into two logs.
    have h_decomp : Filter.Tendsto (fun N => -∑ n ∈ Finset.range N, if n = 0 then 0 else Real.cos (n * (theta + Real.pi / 6)) / n) Filter.atTop (nhds (Real.log (2 * Real.sin (theta / 2 + Real.pi / 12)))) ∧ Filter.Tendsto (fun N => -∑ n ∈ Finset.range N, if n = 0 then 0 else Real.cos (n * (theta + 5 * Real.pi / 6)) / n) Filter.atTop (nhds (Real.log (2 * Real.sin (theta / 2 + 5 * Real.pi / 12)))) := by
      constructor;
      · convert log_two_sin_fourier_series_limit ( theta / 2 + Real.pi / 12 ) ⟨ _, _ ⟩ using 2 <;> ring_nf at * <;> norm_num at * <;> try linarith [ Real.pi_pos, htheta.1, htheta.2 ] ;
        ac_rfl;
      · convert log_two_sin_fourier_series_limit ( theta / 2 + 5 * Real.pi / 12 ) _ using 1;
        · grind;
        · constructor <;> linarith [ htheta.1, htheta.2, Real.pi_pos ];
    convert h_decomp.1.add h_decomp.2 using 2 <;> norm_num [ add_div, Finset.sum_add_distrib ];
    · rw [ ← neg_add, ← Finset.sum_add_distrib ] ; congr ; ext ; aesop;
    · exact?

/-
The integral of theta * cos(n(theta + pi/2)) from 0 to pi/6 is given by the formula.
-/
lemma integral_theta_cos_n_theta_plus_pi_div_two (n : ℕ) (hn : n ≠ 0) :
  ∫ theta in (0 : ℝ)..Real.pi / 6, theta * Real.cos (n * (theta + Real.pi / 2)) =
  (Real.pi / (6 * n)) * Real.sin (2 * n * Real.pi / 3) + (Real.cos (2 * n * Real.pi / 3) - Real.cos (n * Real.pi / 2)) / n ^ 2 := by
    rw [ intervalIntegral.integral_deriv_eq_sub' ];
    case f => exact fun x => x / n * Real.sin ( n * ( x + Real.pi / 2 ) ) + ( Real.cos ( n * ( x + Real.pi / 2 ) ) - Real.cos ( n * ( 0 + Real.pi / 2 ) ) ) / n ^ 2;
    · grind;
    · ext; norm_num [ hn, mul_comm ] ; ring;
      simp +decide [ sq, mul_assoc, hn ];
    · norm_num [ mul_comm ];
    · fun_prop (disch := norm_num)

/-
The partial sums of cos(n(theta + pi/6)) are uniformly bounded.
-/
lemma partial_sums_cos_bounded_1 :
  ∃ C, ∀ N, ∀ theta ∈ Set.Icc 0 (Real.pi / 6), |∑ n ∈ Finset.range N, Real.cos (n * (theta + Real.pi / 6))| ≤ C := by
    -- The sum $\sum_{n=0}^{N-1} \cos(nx)$ is bounded by $1/|\sin(x/2)|$.
    have h_cos_sum_bound : ∀ N : ℕ, ∀ x : ℝ, 0 < x ∧ x < Real.pi → |∑ n ∈ Finset.range N, Real.cos (n * x)| ≤ 1 / |Real.sin (x / 2)| := by
      intros N x hx
      have h_cos_sum_bound : |∑ n ∈ Finset.range N, Real.cos (n * x)| ≤ 1 / |Real.sin (x / 2)| := by
        have h_cos_sum : ∀ N : ℕ, ∑ n ∈ Finset.range N, Real.cos (n * x) = (Real.sin (N * x / 2) * Real.cos ((N - 1) * x / 2)) / Real.sin (x / 2) := by
          intro N; rw [ eq_div_iff ( ne_of_gt ( Real.sin_pos_of_pos_of_lt_pi ( by linarith ) ( by linarith ) ) ) ] ; induction N <;> simp_all +decide [ Finset.sum_range_succ, add_mul ] ; ring; (
          norm_num [ Real.sin_add, Real.cos_add ] ; ring; norm_num [ mul_div ] ; ring;
          rw [ Real.sin_sq, Real.cos_sq ] ; ring;);
        rw [ h_cos_sum, abs_div ];
        gcongr ; exact abs_le.mpr ⟨ by nlinarith [ Real.sin_sq_add_cos_sq ( N * x / 2 ), Real.sin_sq_add_cos_sq ( ( N - 1 ) * x / 2 ) ], by nlinarith [ Real.sin_sq_add_cos_sq ( N * x / 2 ), Real.sin_sq_add_cos_sq ( ( N - 1 ) * x / 2 ) ] ⟩ ;
      exact h_cos_sum_bound;
    use 1 / Real.sin (Real.pi / 12);
    intro N theta htheta; specialize h_cos_sum_bound N ( theta + Real.pi / 6 ) ⟨ by linarith [ Real.pi_pos, htheta.1 ], by linarith [ Real.pi_pos, htheta.2 ] ⟩ ; simp_all +decide [ abs_of_nonneg, Real.sin_nonneg_of_nonneg_of_le_pi ] ;
    exact h_cos_sum_bound.trans ( inv_anti₀ ( Real.sin_pos_of_pos_of_lt_pi ( by positivity ) ( by linarith [ Real.pi_pos ] ) ) ( by rw [ abs_of_nonneg ( Real.sin_nonneg_of_nonneg_of_le_pi ( by linarith [ Real.pi_pos ] ) ( by linarith [ Real.pi_pos ] ) ) ] ; exact Real.sin_le_sin_of_le_of_le_pi_div_two ( by linarith [ Real.pi_pos ] ) ( by linarith [ Real.pi_pos ] ) ( by linarith [ Real.pi_pos ] ) ) )

/-
The partial sums of cos(n(theta + 5pi/6)) are uniformly bounded.
-/
lemma partial_sums_cos_bounded_2 :
  ∃ C, ∀ N, ∀ theta ∈ Set.Icc 0 (Real.pi / 6), |∑ n ∈ Finset.range N, Real.cos (n * (theta + 5 * Real.pi / 6))| ≤ C := by
    -- The sum is bounded by $1/|\sin(x/2)|$ where $x = \theta + 5\pi/6$.
    have h_bound : ∀ N : ℕ, ∀ theta ∈ Set.Icc 0 (Real.pi / 6), |∑ n ∈ Finset.range N, Real.cos (n * (theta + 5 * Real.pi / 6))| ≤ 1 / |Real.sin ((theta + 5 * Real.pi / 6) / 2)| := by
      intros N theta htheta
      have h_sum_cos : ∑ n ∈ Finset.range N, Real.cos (n * (theta + 5 * Real.pi / 6)) = (Real.sin (N * (theta + 5 * Real.pi / 6) / 2) * Real.cos ((N - 1) * (theta + 5 * Real.pi / 6) / 2)) / Real.sin ((theta + 5 * Real.pi / 6) / 2) := by
        rw [ eq_div_iff ];
        · induction N <;> simp_all +decide [ Finset.sum_range_succ, add_mul ];
          rw [ ← Complex.ofReal_inj ] ; norm_num [ Complex.sin, Complex.cos ] ; ring;
          norm_num [ mul_assoc, ← Complex.exp_add ] ; ring;
        · exact ne_of_gt ( Real.sin_pos_of_pos_of_lt_pi ( by linarith [ Real.pi_pos, htheta.1 ] ) ( by linarith [ Real.pi_pos, htheta.2 ] ) );
      rw [ h_sum_cos, abs_div, abs_mul ];
      gcongr ; exact mul_le_one₀ ( Real.abs_sin_le_one _ ) ( abs_nonneg _ ) ( Real.abs_cos_le_one _ ) ;
    -- Since $\sin(x/2) \geq \sin(5\pi/12)$ for $x \in [5\pi/6, \pi]$, we can conclude that $1/|\sin(x/2)| \leq 1/\sin(5\pi/12)$.
    have h_sin_bound : ∀ theta ∈ Set.Icc 0 (Real.pi / 6), |Real.sin ((theta + 5 * Real.pi / 6) / 2)| ≥ Real.sin (5 * Real.pi / 12) := by
      intro theta htheta; rw [ abs_of_nonneg ( Real.sin_nonneg_of_nonneg_of_le_pi ( by linarith [ Set.mem_Icc.mp htheta ] ) ( by linarith [ Set.mem_Icc.mp htheta ] ) ) ] ; rw [ ← Real.cos_pi_div_two_sub ] ; rw [ ← Real.cos_pi_div_two_sub ] ; exact Real.cos_le_cos_of_nonneg_of_le_pi ( by linarith [ Set.mem_Icc.mp htheta ] ) ( by linarith [ Set.mem_Icc.mp htheta ] ) ( by linarith [ Set.mem_Icc.mp htheta ] ) ;
    exact ⟨ 1 / Real.sin ( 5 * Real.pi / 12 ), fun N theta htheta => le_trans ( h_bound N theta htheta ) ( one_div_le_one_div_of_le ( Real.sin_pos_of_pos_of_lt_pi ( by positivity ) ( by linarith [ Real.pi_pos ] ) ) ( h_sin_bound theta htheta ) ) ⟩

/-
The partial sums of the sum of cosines are uniformly bounded.
-/
lemma partial_sums_cos_sum_bounded :
  ∃ C, ∀ N, ∀ theta ∈ Set.Icc 0 (Real.pi / 6),
  |∑ n ∈ Finset.range N, (Real.cos (n * (theta + Real.pi / 6)) + Real.cos (n * (theta + 5 * Real.pi / 6)))| ≤ C := by
    -- Apply the bounds from `partial_sums_cos_bounded_1` and `partial_sums_cos_bounded_2`.
    obtain ⟨C1, hC1⟩ := partial_sums_cos_bounded_1
    obtain ⟨C2, hC2⟩ := partial_sums_cos_bounded_2;
    exact ⟨ C1 + C2, fun N theta htheta => by rw [ Finset.sum_add_distrib ] ; exact abs_le.mpr ⟨ by linarith [ abs_le.mp ( hC1 N theta htheta ), abs_le.mp ( hC2 N theta htheta ) ], by linarith [ abs_le.mp ( hC1 N theta htheta ), abs_le.mp ( hC2 N theta htheta ) ] ⟩ ⟩

/-
If partial sums of a are bounded by C and b is antitone and non-negative, then partial sums of a*b are bounded by 2*C*b(0).
-/
lemma sum_mul_bounded_of_sum_bounded_of_antitone_nonneg {a : ℕ → ℝ} {b : ℕ → ℝ} (C : ℝ)
  (ha : ∀ N, |∑ n ∈ Finset.range N, a n| ≤ C)
  (hb_mono : Antitone b) (hb_nonneg : ∀ n, 0 ≤ b n) :
  ∀ N, |∑ n ∈ Finset.range N, a n * b n| ≤ 2 * C * b 0 := by
    -- By summation by parts, we have $\sum_{n=0}^{N-1} a_n b_n = S_N b_{N-1} + \sum_{n=0}^{N-2} S_{n+1} (b_n - b_{n+1})$.
    have h_summation_by_parts : ∀ N, ∑ n ∈ Finset.range N, a n * b n = (∑ n ∈ Finset.range N, a n) * b (N - 1) + ∑ n ∈ Finset.range (N - 1), (∑ n ∈ Finset.range (n + 1), a n) * (b n - b (n + 1)) := by
      intro N; induction' N with N ih <;> simp_all +decide [ Finset.sum_range_succ ] ; ring;
      cases N <;> norm_num [ Finset.sum_range_succ ] ; ring;
    -- Applying the bounds from ha and hb_mono, we get:
    have h_bound : ∀ N, |(∑ n ∈ Finset.range N, a n) * b (N - 1)| ≤ C * b 0 ∧ |∑ n ∈ Finset.range (N - 1), (∑ n ∈ Finset.range (n + 1), a n) * (b n - b (n + 1))| ≤ C * b 0 := by
      intro N
      constructor;
      · rw [ abs_mul, abs_of_nonneg ( hb_nonneg _ ) ] ; exact mul_le_mul ( ha _ ) ( hb_mono ( Nat.zero_le _ ) ) ( hb_nonneg _ ) ( by linarith [ abs_le.mp ( ha N ) ] ) ;
      · refine' le_trans ( Finset.abs_sum_le_sum_abs _ _ ) _;
        refine' le_trans ( Finset.sum_le_sum fun i hi => _ ) _;
        use fun i => C * ( b i - b ( i + 1 ) );
        · rw [ abs_mul, abs_of_nonneg ( sub_nonneg.mpr ( hb_mono ( Nat.le_succ _ ) ) ) ] ; exact mul_le_mul_of_nonneg_right ( ha _ ) ( sub_nonneg.mpr ( hb_mono ( Nat.le_succ _ ) ) ) ;
        · rw [ ← Finset.mul_sum _ _ _ ];
          exact mul_le_mul_of_nonneg_left ( by rw [ Finset.sum_range_sub' ] ; linarith [ hb_nonneg 0, hb_nonneg ( N - 1 ) ] ) ( show 0 ≤ C by exact le_trans ( abs_nonneg _ ) ( ha 0 ) );
    exact fun N => by rw [ h_summation_by_parts N ] ; exact abs_le.mpr ⟨ by linarith [ abs_le.mp ( h_bound N |>.1 ), abs_le.mp ( h_bound N |>.2 ) ], by linarith [ abs_le.mp ( h_bound N |>.1 ), abs_le.mp ( h_bound N |>.2 ) ] ⟩ ;

/-
The partial sums of the series are uniformly bounded.
-/
lemma partial_sums_series_bounded :
  ∃ C, ∀ N, ∀ theta ∈ Set.Icc 0 (Real.pi / 6),
  |∑ n ∈ Finset.range N, if n = 0 then 0 else (Real.cos (n * (theta + Real.pi / 6)) + Real.cos (n * (theta + 5 * Real.pi / 6))) / n| ≤ C := by
    obtain ⟨ C₁, hC₁ ⟩ := @partial_sums_cos_sum_bounded;
    use 2 * (C₁ + 2);
    intros N theta htheta
    have h_partial_sum : ∀ K : ℕ, |∑ n ∈ Finset.range K, (Real.cos ((n + 1) * (theta + Real.pi / 6)) + Real.cos ((n + 1) * (theta + 5 * Real.pi / 6))) / (n + 1 : ℝ)| ≤ 2 * (C₁ + 2) := by
      intros K
      have h_partial_sum : ∀ K : ℕ, |∑ n ∈ Finset.range K, (Real.cos ((n + 1) * (theta + Real.pi / 6)) + Real.cos ((n + 1) * (theta + 5 * Real.pi / 6)))| ≤ C₁ + 2 := by
        intro K
        specialize hC₁ (K + 1) theta htheta
        have h_partial_sum : |∑ n ∈ Finset.range (K + 1), (Real.cos (n * (theta + Real.pi / 6)) + Real.cos (n * (theta + 5 * Real.pi / 6)))| ≤ C₁ := by
          exact hC₁;
        simp_all +decide [ Finset.sum_range_succ' ];
        exact abs_le.mpr ⟨ by linarith [ abs_le.mp hC₁ ], by linarith [ abs_le.mp hC₁ ] ⟩;
      convert sum_mul_bounded_of_sum_bounded_of_antitone_nonneg ( C₁ + 2 ) ( fun K => h_partial_sum K ) ( show Antitone ( fun n : ℕ => ( n + 1 : ℝ ) ⁻¹ ) from fun n m hnm => inv_anti₀ ( by positivity ) ( by norm_cast; linarith ) ) ( fun n => by positivity ) K using 1;
      norm_num;
    convert h_partial_sum ( N - 1 ) using 1 ; rcases N <;> norm_num [ Finset.sum_range_succ' ] at *

/-
The partial sums of the shifted cosine terms are uniformly bounded.
-/
lemma partial_sums_shifted_cos_bounded :
  ∃ C, ∀ N, ∀ theta ∈ Set.Icc 0 (Real.pi / 6),
  |∑ n ∈ Finset.range N, (Real.cos ((n + 1) * (theta + Real.pi / 6)) + Real.cos ((n + 1) * (theta + 5 * Real.pi / 6)))| ≤ C := by
    -- By the properties of the cosine function and the periodicity, we can find such a constant $C$.
    have h_bound : ∃ C : ℝ, ∀ N : ℕ, ∀ theta ∈ Set.Icc 0 (Real.pi / 6), |∑ n ∈ Finset.range N, (Real.cos (n * (theta + Real.pi / 6)) + Real.cos (n * (theta + 5 * Real.pi / 6)))| ≤ C := by
      exact?
    obtain ⟨C, hC⟩ := h_bound;
    use C + 2; intros N theta htheta; (
    have := hC ( N + 1 ) theta htheta; simp_all +decide [ Finset.sum_range_succ' ] ;
    exact abs_le.mpr ⟨ by linarith [ abs_le.mp this ], by linarith [ abs_le.mp this ] ⟩);

/-
Definition of S_series_term.
-/
noncomputable def S_series_term (n : ℕ) : ℝ :=
  if n = 0 then 0 else
  -8 * (Real.cos (n * Real.pi / 3) / n) * ∫ theta in (0 : ℝ)..Real.pi / 6, theta * Real.cos (n * (theta + Real.pi / 2))

/-
S is equal to the sum of S_series_term.
-/
lemma S_eq_series_sum : S_sum = ∑' n : ℕ, S_series_term n := by
  -- By definition of $S_sum$, we know that
  have hS_sum_def : S_sum = ∫ theta in (0 : ℝ)..Real.pi / 6, 4 * theta * Real.log (1 + 2 * Real.sin theta) := by
    rw [ S_eq_S_integral_theta ];
    unfold S_integral_theta;
    rw [ intervalIntegral.integral_of_le ( by positivity ), intervalIntegral.integral_of_le ( by positivity ) ];
    rw [ MeasureTheory.integral_Ioc_eq_integral_Ioo, MeasureTheory.integral_Ioc_eq_integral_Ioo ];
    refine' MeasureTheory.setIntegral_congr_fun measurableSet_Ioo fun x hx => _;
    convert integrand_simplification x ⟨ hx.1, hx.2 ⟩ using 1;
  -- Using the series expansion of $\log(1 + 2\sin\theta)$, we can rewrite the integral.
  have h_series : ∀ N, ∫ theta in (0 : ℝ)..Real.pi / 6, 4 * theta * (-∑ n ∈ Finset.range N, if n = 0 then 0 else (Real.cos (n * (theta + Real.pi / 6)) + Real.cos (n * (theta + 5 * Real.pi / 6))) / n) = ∑ n ∈ Finset.range N, S_series_term n := by
    intro N
    simp [S_series_term];
    rw [ intervalIntegral.integral_congr fun x hx => by rw [ Finset.mul_sum _ _ _ ] ] ; rw [ intervalIntegral.integral_finset_sum ] ; norm_num [ mul_assoc, mul_comm, mul_left_comm, div_eq_mul_inv ] ; ring;
    · rw [ ← Finset.sum_neg_distrib ] ; refine' Finset.sum_congr rfl fun n hn => _ ; by_cases hn' : n = 0 <;> simp +decide [ hn', mul_assoc, mul_comm, mul_left_comm, div_eq_mul_inv ] ; ring;
      rw [ intervalIntegral.integral_add ] <;> norm_num ; ring;
      · norm_num [ mul_assoc, mul_comm, mul_left_comm, ← intervalIntegral.integral_const_mul ] ; ring;
        rw [ ← intervalIntegral.integral_add ] ; congr ; ext x ; ring;
        · rw [ show x * n + n * Real.pi * ( 1 / 6 ) = x * n + n * Real.pi * ( 1 / 2 ) - n * Real.pi * ( 1 / 3 ) by ring, show x * n + n * Real.pi * ( 5 / 6 ) = x * n + n * Real.pi * ( 1 / 2 ) + n * Real.pi * ( 1 / 3 ) by ring ] ; norm_num [ Real.cos_sub, Real.cos_add ] ; ring;
        · exact Continuous.intervalIntegrable ( by exact Continuous.mul ( Continuous.mul ( continuous_id.mul continuous_const ) ( Real.continuous_cos.comp ( by continuity ) ) ) continuous_const ) _ _;
        · exact Continuous.intervalIntegrable ( by exact Continuous.mul ( Continuous.mul ( Continuous.mul continuous_const continuous_id' ) ( Real.continuous_cos.comp <| by continuity ) ) continuous_const ) _ _;
      · exact Continuous.intervalIntegrable ( by exact Continuous.mul ( Continuous.mul ( continuous_id.mul continuous_const ) ( Real.continuous_cos.comp ( by continuity ) ) ) continuous_const ) _ _;
      · exact Continuous.intervalIntegrable ( by exact Continuous.mul ( Continuous.mul ( continuous_id.mul continuous_const ) ( Real.continuous_cos.comp ( by continuity ) ) ) continuous_const ) _ _;
    · exact fun i hi => Continuous.intervalIntegrable ( by split_ifs <;> continuity ) _ _;
  -- By the Dominated Convergence Theorem, we can interchange the limit and the integral.
  have h_dominated_convergence : Filter.Tendsto (fun N => ∫ theta in (0 : ℝ)..Real.pi / 6, 4 * theta * (-∑ n ∈ Finset.range N, if n = 0 then 0 else (Real.cos (n * (theta + Real.pi / 6)) + Real.cos (n * (theta + 5 * Real.pi / 6))) / n)) Filter.atTop (nhds (∫ theta in (0 : ℝ)..Real.pi / 6, 4 * theta * Real.log (1 + 2 * Real.sin theta))) := by
    -- The partial sums of the series are uniformly bounded and converge pointwise to the logarithm function.
    have h_uniform_bounded : ∃ C, ∀ N, ∀ theta ∈ Set.Icc 0 (Real.pi / 6), |∑ n ∈ Finset.range N, if n = 0 then 0 else (Real.cos (n * (theta + Real.pi / 6)) + Real.cos (n * (theta + 5 * Real.pi / 6))) / n| ≤ C := by
      exact?;
    refine' intervalIntegral.tendsto_integral_filter_of_dominated_convergence _ _ _ _ _;
    use fun x => 4 * |x| * h_uniform_bounded.choose;
    · refine' Filter.Eventually.of_forall fun N => Continuous.aestronglyMeasurable _;
      refine' Continuous.mul ( continuous_const.mul continuous_id' ) _;
      exact Continuous.neg <| continuous_finset_sum _ fun _ _ => by split_ifs <;> [ exact continuous_const; exact Continuous.div_const ( Continuous.add ( Real.continuous_cos.comp <| by continuity ) ( Real.continuous_cos.comp <| by continuity ) ) _ ] ;
    · filter_upwards [ Filter.eventually_gt_atTop 0 ] with N hN;
      filter_upwards [ ] with x hx using by simpa [ abs_mul ] using mul_le_mul_of_nonneg_left ( h_uniform_bounded.choose_spec N x <| by constructor <;> cases Set.mem_uIoc.mp hx <;> linarith [ Real.pi_pos ] ) ( by positivity ) ;
    · exact Continuous.intervalIntegrable ( by continuity ) _ _;
    · refine' Filter.Eventually.of_forall fun x hx => _;
      have := log_one_plus_two_sin_series_limit x ⟨ by cases Set.mem_uIoc.mp hx <;> linarith [ Real.pi_pos ], by cases Set.mem_uIoc.mp hx <;> linarith [ Real.pi_pos ] ⟩;
      exact this.const_mul _;
  refine' tendsto_nhds_unique h_dominated_convergence _ |> fun h => hS_sum_def.trans h |> fun h => h.trans _;
  convert Summable.hasSum _ |> HasSum.tendsto_sum_nat;
  exact h_series _;
  · have h_abs_conv : Summable (fun n : ℕ => |S_series_term n|) := by
      have h_abs_conv : Summable (fun n : ℕ => |Real.cos (n * Real.pi / 3) / n * ∫ theta in (0 : ℝ)..Real.pi / 6, theta * Real.cos (n * (theta + Real.pi / 2))|) := by
        have h_abs_conv : Summable (fun n : ℕ => |Real.cos (n * Real.pi / 3) / n * ((Real.pi / (6 * n)) * Real.sin (2 * n * Real.pi / 3) + (Real.cos (2 * n * Real.pi / 3) - Real.cos (n * Real.pi / 2)) / n ^ 2)|) := by
          -- We'll use the fact that |cos(x)| ≤ 1 and |sin(x)| ≤ 1 for all x.
          have h_bound : ∀ n : ℕ, |Real.cos (n * Real.pi / 3) / n * ((Real.pi / (6 * n)) * Real.sin (2 * n * Real.pi / 3) + (Real.cos (2 * n * Real.pi / 3) - Real.cos (n * Real.pi / 2)) / n ^ 2)| ≤ (1 / n) * ((Real.pi / (6 * n)) + (2 / n ^ 2)) := by
            intro n
            simp [abs_mul];
            refine' mul_le_mul _ _ _ _ <;> norm_num [ abs_div, abs_mul, abs_of_nonneg, Real.pi_pos.le ];
            · exact mul_le_of_le_one_left ( by positivity ) ( Real.abs_cos_le_one _ ) |> le_trans <| by norm_num;
            · refine' abs_le.mpr ⟨ _, _ ⟩ <;> ring_nf <;> norm_num [ Real.pi_pos.le ];
              · nlinarith [ Real.pi_pos, show ( 0 : ℝ ) ≤ Real.pi * ( n : ℝ ) ⁻¹ by positivity, show ( 0 : ℝ ) ≤ ( n ^ 2 : ℝ ) ⁻¹ by positivity, abs_le.mp ( Real.abs_sin_le_one ( Real.pi * n * ( 2 / 3 ) ) ), abs_le.mp ( Real.abs_cos_le_one ( Real.pi * n * ( 2 / 3 ) ) ), abs_le.mp ( Real.abs_cos_le_one ( Real.pi * n * ( 1 / 2 ) ) ) ];
              · nlinarith [ abs_le.mp ( Real.abs_sin_le_one ( Real.pi * n * ( 2 / 3 ) ) ), abs_le.mp ( Real.abs_cos_le_one ( Real.pi * n * ( 2 / 3 ) ) ), abs_le.mp ( Real.abs_cos_le_one ( Real.pi * n * ( 1 / 2 ) ) ), show 0 ≤ Real.pi * ( n : ℝ ) ⁻¹ by positivity, show 0 ≤ ( n ^ 2 : ℝ ) ⁻¹ by positivity ];
          refine' Summable.of_nonneg_of_le ( fun n => abs_nonneg _ ) ( fun n => h_bound n ) _;
          ring_nf;
          exact Summable.add ( Summable.mul_right _ <| Summable.mul_right _ <| by simp ) ( Summable.mul_right _ <| by simp );
        refine' h_abs_conv.congr fun n => _;
        by_cases hn : n = 0 <;> simp +decide [ hn, integral_theta_cos_n_theta_plus_pi_div_two ]
      convert h_abs_conv.mul_left 8 using 2 ; norm_num [ S_series_term ] ; ring;
      split_ifs <;> norm_num [ abs_mul, abs_neg ] ; ring;
      grind;
    exact h_abs_conv.of_abs;
  · rfl

/-
Definition of a_n.
-/
noncomputable def S_a_n (n : ℕ) : ℝ := Real.cos (n * Real.pi / 3) * (Real.cos (2 * n * Real.pi / 3) - Real.cos (n * Real.pi / 2))

/-
Definition of C.
-/
noncomputable def C_sum : ℝ := ∑' n : ℕ, if n = 0 then 0 else Real.sin (n * Real.pi / 3) / n^2

/-
cos(n*pi/3) * sin(2*n*pi/3) = 1/2 * sin(n*pi/3).
-/
lemma cos_mul_sin_identity (n : ℕ) : Real.cos (n * Real.pi / 3) * Real.sin (2 * n * Real.pi / 3) = 1 / 2 * Real.sin (n * Real.pi / 3) := by
  rw [ show 2 * n * Real.pi / 3 = n * Real.pi - n * Real.pi / 3 by ring ] ; rw [ Real.sin_sub ] ; ring ; norm_num;
  rcases Nat.even_or_odd' n with ⟨ k, rfl | rfl ⟩ <;> norm_num <;> ring_nf <;> norm_num [ mul_div, mul_assoc, mul_comm Real.pi ] at *;
  · rw [ ← Nat.mod_add_div k 3 ] ; norm_num [ mul_add, add_mul, mul_assoc, mul_comm Real.pi _, mul_div_assoc ] ; ring_nf; norm_num [ Nat.add_mod, Nat.mul_mod ] ;
    have := Nat.mod_lt k zero_lt_three; interval_cases k % 3 <;> norm_num [ mul_assoc, mul_comm Real.pi ] ; ring; norm_num;
    · norm_num [ mul_two, Real.sin_add ];
    · norm_num [ show 2 / 3 * Real.pi = Real.pi - Real.pi / 3 by ring ] ; ring;
    · norm_num [ ( by ring : 2 * ( 2 / 3 * Real.pi ) = Real.pi + Real.pi / 3 ), Real.sin_add, Real.cos_add ] ; ring;
  · rw [ ← Nat.mod_add_div k 3 ] ; have := Nat.mod_lt k zero_lt_three; interval_cases k % 3 <;> norm_num [ Real.cos_add, Real.sin_add, mul_div_assoc ] ; ring;
    · norm_num [ mul_two, Real.sin_add, Real.cos_add, mul_div ] ; ring;
      norm_num [ show Real.cos ( ( k / 3 : ℕ ) * Real.pi ) ^ 4 = ( Real.cos ( ( k / 3 : ℕ ) * Real.pi ) ^ 2 ) ^ 2 by ring, Real.cos_sq' ];
    · norm_num [ ( by ring : 1 / 3 * Real.pi = Real.pi / 3 ), Real.cos_add, Real.sin_add, mul_div_assoc ] ; ring; norm_num;
      norm_num [ mul_two, Real.sin_add, Real.cos_add, show Real.pi * ( 2 / 3 ) = Real.pi - Real.pi / 3 by ring ] ; ring; norm_num; ring;
      norm_num [ pow_three ] ; ring;
    · ring_nf; norm_num [ mul_div, mul_assoc, mul_comm Real.pi _, mul_left_comm ] ; ring;
      norm_num [ mul_assoc, mul_comm Real.pi _, mul_div ] ; ring; norm_num [ Real.sin_add, Real.cos_add, mul_div ] ; ring;
      rw [ show Real.pi * ( 4 / 3 ) = Real.pi + Real.pi / 3 by ring, Real.cos_add, Real.sin_add ] ; norm_num ; ring;
      norm_num [ pow_three ] ; ring

/-
S = -2pi/3 * C - 8 * sum(a_n/n^3).
-/
lemma S_eq_final_form : S_sum = - (2 * Real.pi / 3) * C_sum - 8 * (∑' n : ℕ, if n = 0 then 0 else S_a_n n / n^3) := by
  have h_simp : ∀ n : ℕ, S_series_term n = -8 * (Real.cos (n * Real.pi / 3) / n) * ((Real.pi / (6 * n)) * Real.sin (2 * n * Real.pi / 3) + (Real.cos (2 * n * Real.pi / 3) - Real.cos (n * Real.pi / 2)) / n ^ 2) := by
    intro n
    unfold S_series_term
    by_cases hn : n = 0 <;> simp [hn];
    exact Or.inl <| integral_theta_cos_n_theta_plus_pi_div_two n hn ▸ by ring;
  have h_series_sum : ∑' n : ℕ, S_series_term n = -8 * (∑' n : ℕ, if n = 0 then 0 else (Real.cos (n * Real.pi / 3) / n) * ((Real.pi / (6 * n)) * Real.sin (2 * n * Real.pi / 3))) - 8 * (∑' n : ℕ, if n = 0 then 0 else (Real.cos (n * Real.pi / 3) / n) * ((Real.cos (2 * n * Real.pi / 3) - Real.cos (n * Real.pi / 2)) / n ^ 2)) := by
    have h_series_sum : Summable (fun n : ℕ => if n = 0 then 0 else (Real.cos (n * Real.pi / 3) / n) * ((Real.pi / (6 * n)) * Real.sin (2 * n * Real.pi / 3))) ∧ Summable (fun n : ℕ => if n = 0 then 0 else (Real.cos (n * Real.pi / 3) / n) * ((Real.cos (2 * n * Real.pi / 3) - Real.cos (n * Real.pi / 2)) / n ^ 2)) := by
      constructor <;> norm_num [ div_eq_mul_inv, mul_assoc, mul_comm, mul_left_comm ];
      · -- The series $\sum_{n=1}^{\infty} \frac{\cos(n\pi/3) \sin(2n\pi/3)}{n^2}$ is absolutely convergent.
        have h_abs_conv : Summable (fun n : ℕ => |Real.cos (n * Real.pi / 3) * Real.sin (2 * n * Real.pi / 3) / n ^ 2|) := by
          exact Summable.of_nonneg_of_le ( fun n => abs_nonneg _ ) ( fun n => by simpa [ abs_div, abs_mul ] using div_le_div_of_nonneg_right ( mul_le_mul ( Real.abs_cos_le_one _ ) ( Real.abs_sin_le_one _ ) ( by positivity ) ( by positivity ) ) ( sq_nonneg _ ) ) ( Real.summable_one_div_nat_pow.2 one_lt_two );
        convert h_abs_conv.of_abs.mul_left ( Real.pi / 6 ) using 2 ; ring ; aesop;
      · field_simp;
        -- The absolute value of the terms is bounded by $4/n^3$, and since $\sum' n : ℕ, 4/n^3$ converges, the original series is absolutely convergent.
        have h_abs : ∀ n : ℕ, |if n = 0 then 0 else Real.cos (n * Real.pi / 3) * (Real.cos (n * Real.pi * 2 / 3) - Real.cos (n * Real.pi / 2)) / n^3| ≤ 4 / n^3 := by
          intro n; split_ifs <;> norm_num [ abs_div, abs_mul ];
          · positivity;
          · gcongr ; exact le_trans ( mul_le_mul ( Real.abs_cos_le_one _ ) ( abs_sub _ _ ) ( by positivity ) ( by positivity ) ) ( by linarith [ Real.abs_cos_le_one ( n * Real.pi * 2 / 3 ), Real.abs_cos_le_one ( n * Real.pi / 2 ) ] ) ;
        -- Since the absolute value of the terms is bounded by a convergent p-series, the original series is absolutely convergent.
        have h_abs_conv : Summable (fun n : ℕ => |if n = 0 then 0 else Real.cos (n * Real.pi / 3) * (Real.cos (n * Real.pi * 2 / 3) - Real.cos (n * Real.pi / 2)) / n^3|) := by
          exact Summable.of_nonneg_of_le ( fun n => abs_nonneg _ ) h_abs ( Summable.mul_left _ <| Real.summable_nat_pow_inv.2 <| by norm_num );
        exact h_abs_conv.of_abs;
    simp_all +decide [ mul_add, Finset.sum_add_distrib, tsum_mul_left, tsum_mul_right ];
    rw [ Summable.tsum_add ];
    · rw [ tsum_neg, tsum_neg ] ; norm_num [ mul_assoc, tsum_mul_left ] ; ring;
      exact congrArg₂ _ ( by congr; ext n; by_cases hn : n = 0 <;> simp +decide [ hn ] ; ring ) ( by congr; ext n; by_cases hn : n = 0 <;> simp +decide [ hn ] );
    · convert h_series_sum.1.mul_left ( -8 ) using 2 ; ring;
      split_ifs <;> ring ; aesop;
    · convert h_series_sum.2.mul_left ( -8 ) using 2 ; ring;
      split_ifs <;> ring ; aesop;
  convert h_series_sum using 2;
  · exact?;
  · have h_cos_sin : ∀ n : ℕ, Real.cos (n * Real.pi / 3) * Real.sin (2 * n * Real.pi / 3) = 1 / 2 * Real.sin (n * Real.pi / 3) := by
      exact?;
    have h_series_sum : ∑' n : ℕ, (if n = 0 then 0 else (Real.cos (n * Real.pi / 3) / n) * ((Real.pi / (6 * n)) * Real.sin (2 * n * Real.pi / 3))) = (Real.pi / 12) * ∑' n : ℕ, (if n = 0 then 0 else Real.sin (n * Real.pi / 3) / n ^ 2) := by
      rw [ ← tsum_mul_left ] ; congr ; ext n ; by_cases hn : n = 0 <;> simp +decide [ hn, h_cos_sin ] ; ring;
      grind;
    rw [ h_series_sum ] ; ring!;
  · unfold S_a_n; congr; ext n; ring;

/-
Values of a_n for n = 0 to 11.
-/
lemma S_a_n_values (n : ℕ) (hn : n < 12) :
  S_a_n n = [0, -1/4, -1/4, -1, 3/4, -1/4, 2, -1/4, 3/4, -1, -1/4, -1/4].get ⟨n, by simp [hn]⟩ := by
    interval_cases n <;> norm_num [ S_a_n ];
    all_goals norm_num [ show 2 * Real.pi / 3 = Real.pi - Real.pi / 3 by ring, show 4 * Real.pi / 3 = Real.pi + Real.pi / 3 by ring, show 6 * Real.pi / 3 = 2 * Real.pi by ring, show 8 * Real.pi / 3 = 2 * Real.pi + 2 * Real.pi / 3 by ring, show 10 * Real.pi / 3 = 3 * Real.pi + Real.pi / 3 by ring, show 12 * Real.pi / 3 = 4 * Real.pi by ring, Real.cos_add, Real.cos_sub, Real.cos_two_mul, Real.cos_three_mul ] at *;
    all_goals norm_num [ show 3 * Real.pi / 2 = Real.pi + Real.pi / 2 by ring, show 4 * Real.pi / 2 = 2 * Real.pi by ring, show 5 * Real.pi / 3 = 2 * Real.pi - Real.pi / 3 by ring, show 6 * Real.pi / 2 = 3 * Real.pi by ring, show 7 * Real.pi / 3 = 2 * Real.pi + Real.pi / 3 by ring, show 8 * Real.pi / 2 = 4 * Real.pi by ring, Real.cos_add, Real.cos_sub ] at *;
    all_goals norm_num [ show 3 * Real.pi = 2 * Real.pi + Real.pi by ring, show 4 * Real.pi = 2 * Real.pi + 2 * Real.pi by ring, show 5 * Real.pi / 2 = 2 * Real.pi + Real.pi / 2 by ring, show 6 * Real.pi / 2 = 3 * Real.pi by ring, show 7 * Real.pi / 2 = 3 * Real.pi + Real.pi / 2 by ring, show 8 * Real.pi / 2 = 4 * Real.pi by ring, Real.sin_add, Real.cos_add ] at *;
    · norm_num [ show 14 * Real.pi / 3 = 4 * Real.pi + 2 * Real.pi / 3 by ring, Real.cos_add, show 4 * Real.pi = 2 * Real.pi + 2 * Real.pi by ring, Real.cos_two_mul, mul_div_assoc ];
    · norm_num [ show 16 * Real.pi / 3 = 2 * ( 8 * Real.pi / 3 ) by ring, Real.cos_two_mul, show 8 * Real.pi / 3 = 2 * ( 4 * Real.pi / 3 ) by ring, Real.cos_two_mul, show 4 * Real.pi / 3 = Real.pi + Real.pi / 3 by ring, Real.cos_add ];
    · norm_num [ show 9 * Real.pi / 3 = 3 * Real.pi by ring, show 18 * Real.pi / 3 = 6 * Real.pi by ring, show 9 * Real.pi / 2 = 4 * Real.pi + Real.pi / 2 by ring, Real.cos_three_mul, Real.cos_two_mul, Real.cos_add ];
      norm_num [ show 4 * Real.pi = 2 * Real.pi + 2 * Real.pi by ring, show 6 * Real.pi = 4 * Real.pi + 2 * Real.pi by ring ];
    · rw [ show 20 * Real.pi / 3 = 2 * ( 10 * Real.pi / 3 ) by ring, show 10 * Real.pi / 2 = 5 * Real.pi by ring ] ; norm_num [ Real.cos_two_mul, Real.cos_three_mul, show 5 * Real.pi = Real.pi + Real.pi + Real.pi + Real.pi + Real.pi by ring ] ; ring; norm_num;
      norm_num [ ( by ring : Real.pi * ( 10 / 3 ) = Real.pi / 3 + Real.pi + Real.pi + Real.pi ) ];
    · norm_num [ show 11 * Real.pi / 3 = 4 * Real.pi - Real.pi / 3 by ring, show 22 * Real.pi / 3 = 7 * Real.pi + Real.pi / 3 by ring, show 11 * Real.pi / 2 = 5 * Real.pi + Real.pi / 2 by ring, Real.cos_add, Real.cos_sub ] at *;
      norm_num [ show 4 * Real.pi = 2 * Real.pi + 2 * Real.pi by ring, show 7 * Real.pi = 2 * Real.pi + 2 * Real.pi + 2 * Real.pi + Real.pi by ring, show 5 * Real.pi = 2 * Real.pi + 2 * Real.pi + Real.pi by ring ]

/-
a_n is periodic with period 12.
-/
lemma S_a_n_periodic (n : ℕ) : S_a_n (n + 12) = S_a_n n := by
  unfold S_a_n; norm_num [ show 12 * Real.pi / 3 = 4 * Real.pi by ring, show 2 * Real.pi = 2 * Real.pi by ring, Real.cos_add, Real.sin_add, show 4 * Real.pi = 2 * Real.pi + 2 * Real.pi by ring ] ; ring;
  norm_num [ show Real.pi * 4 = 2 * Real.pi + 2 * Real.pi by ring, show Real.pi * 8 = 2 * Real.pi + 2 * Real.pi + 2 * Real.pi + 2 * Real.pi by ring, show Real.pi * 6 = 2 * Real.pi + 2 * Real.pi + 2 * Real.pi by ring, Real.cos_add ]

/-
Formula for a_n using modular arithmetic.
-/
lemma S_a_n_formula (n : ℕ) :
  S_a_n n = -1/4 + (if n % 3 = 0 then -3/4 else 0) + (if n % 4 = 0 then 1 else 0) + (if n % 6 = 0 then 3 else 0) + (if n % 12 = 0 then -3 else 0) := by
    -- By definition of periodicity, we can reduce $n$ modulo 12 and use the values from the table.
    have h_mod : S_a_n n = S_a_n (n % 12) := by
      rw [ ← Nat.mod_add_div n 12 ];
      induction n / 12 <;> simp_all +decide [ Nat.mul_succ, ← add_assoc, S_a_n_periodic ];
    -- By checking each possible value of $n \mod 12$, we can verify that the expression holds.
    have h_cases : ∀ m < 12, S_a_n m = (((-1 / 4 + if m % 3 = 0 then -3 / 4 else 0) + if m % 4 = 0 then 1 else 0) + if m % 6 = 0 then 3 else 0) + if m % 12 = 0 then -3 else 0 := by
      intro m hm
      have h_eq : S_a_n m = [0, -1/4, -1/4, -1, 3/4, -1/4, 2, -1/4, 3/4, -1, -1/4, -1/4].get! m := by
        have := S_a_n_values m hm; aesop;
      interval_cases m <;> norm_num [ h_eq ];
    rw [ h_mod, h_cases _ ( Nat.mod_lt _ ( by decide ) ) ] ; norm_num [ Nat.add_mod, Nat.mul_mod ] ;

/-
Sum of 1/n^s where k divides n is k^(-s) * zeta(s).
-/
lemma sum_indicator_div_pow (k : ℕ) (hk : k ≠ 0) (s : ℂ) (hs : 1 < s.re) :
  ∑' n : ℕ, (if n = 0 then 0 else (if n % k = 0 then 1 else 0) / (n : ℂ)^s) = k^(-s) * riemannZeta s := by
    -- The sum over n divisible by k can be rewritten by substituting n = k * m, where m is a positive integer.
    have h_subst : ∑' n : ℕ, (if n = 0 then 0 else (if n % k = 0 then 1 else 0) / (n : ℂ) ^ s) = ∑' m : ℕ, (if m = 0 then 0 else 1 / ((k * m) : ℂ) ^ s) := by
      have h_subst : {n : ℕ | n ≠ 0 ∧ n % k = 0} = Set.image (fun m => k * m) {m : ℕ | m ≠ 0} := by
        ext n
        simp [Set.mem_image];
        exact ⟨ fun h => ⟨ n / k - 1, by nlinarith [ Nat.mod_add_div n k, Nat.sub_add_cancel ( show 1 ≤ n / k from Nat.div_pos ( Nat.le_of_dvd ( Nat.pos_of_ne_zero h.1 ) ( Nat.dvd_of_mod_eq_zero h.2 ) ) ( Nat.pos_of_ne_zero hk ) ) ] ⟩, by rintro ⟨ m, rfl ⟩ ; exact ⟨ by aesop, by norm_num ⟩ ⟩;
      rw [ Set.ext_iff ] at h_subst;
      rw [ tsum_eq_tsum_of_ne_zero_bij ];
      use fun x => k * x;
      · aesop_cat;
      · intro n hn; specialize h_subst n; aesop;
      · aesop;
    simp_all +decide [ Complex.cpow_neg ];
    rw [ zeta_eq_tsum_one_div_nat_cpow ];
    · rw [ ← tsum_mul_left ] ; congr ; ext n ; by_cases hn : n = 0 <;> simp +decide [ *, mul_comm, Complex.cpow_neg, Complex.cpow_mul ] ;
      · rintro rfl; norm_num at hs;
      · rw [ ← mul_inv, Complex.cpow_def_of_ne_zero, Complex.cpow_def_of_ne_zero, Complex.cpow_def_of_ne_zero ] <;> norm_num [ hn, hk ];
        rw [ ← mul_inv, ← Complex.exp_add, Complex.log_mul ( by norm_cast ) ( by norm_cast ) ] ; ring;
        norm_num [ Complex.arg ];
        exact ⟨ Real.pi_pos, Real.pi_pos.le ⟩;
    · linarith

/-
Trigamma duplication formula.
-/
lemma trigamma_duplication (z : ℝ) : trigamma z + trigamma (z + 1/2) = 4 * trigamma (2 * z) := by
  by_cases hz : z = 0 <;> simp_all +decide [ trigamma ];
  · -- Evaluate the sum $\sum_{n=1}^{\infty} \frac{1}{(n + \frac{1}{2})^2}$.
    have h_sum_half : ∑' n : ℕ, (1 / (n + 1 / 2 : ℝ)^2) = 4 * (∑' n : ℕ, (1 / (2 * n + 1 : ℝ)^2)) := by
      field_simp;
      norm_num [ div_eq_mul_inv, mul_comm, ← tsum_mul_left ];
    -- We can split the sum into two separate sums, one over even terms and one over odd terms.
    have h_split : ∑' n : ℕ, (1 / (n : ℝ) ^ 2) = ∑' n : ℕ, (1 / (2 * n : ℝ) ^ 2) + ∑' n : ℕ, (1 / (2 * n + 1 : ℝ) ^ 2) := by
      rw [ ← tsum_even_add_odd ] <;> norm_num [ mul_pow ];
      · exact Summable.mul_right _ <| Real.summable_nat_pow_inv.2 one_lt_two;
      · exact_mod_cast Summable.comp_injective ( Real.summable_nat_pow_inv.2 one_lt_two ) fun x y h => by simpa using h;
    norm_num [ mul_pow, tsum_mul_right ] at *; have := hasSum_zeta_two; have := this.tsum_eq; norm_num at *; linarith;
  · by_cases h : Summable ( fun n : ℕ => ( ( n : ℝ ) + z ) ⁻¹ ^ 2 ) <;> simp_all +decide [ tsum_eq_zero_of_not_summable ];
    · have h_split : ∑' n : ℕ, (1 / (n + 2 * z : ℝ)^2) = (1 / 4) * (∑' n : ℕ, (1 / (n + z : ℝ)^2) + ∑' n : ℕ, (1 / (n + z + 1 / 2 : ℝ)^2)) := by
        have h_split : ∑' n : ℕ, (1 / (n + 2 * z : ℝ)^2) = (1 / 4) * (∑' n : ℕ, (1 / (n / 2 + z : ℝ)^2)) := by
          rw [ ← tsum_mul_left ] ; congr ; ext n ; rw [ div_mul_div_comm ] ; ring;
        rw [ h_split, ← tsum_even_add_odd ];
        · norm_num [ add_assoc, mul_div_cancel_left₀ ] ; congr ; ext n ; ring;
        · simpa [ mul_div_cancel_left₀ ] using h;
        · have h_summable : Summable (fun k : ℕ => (1 : ℝ) / (k + z + 1 / 2) ^ 2) := by
            rw [ ← summable_nat_add_iff ( ⌈|z|⌉₊ + 1 ) ] at *;
            norm_num +zetaDelta at *;
            exact Summable.of_nonneg_of_le ( fun n => inv_nonneg.2 ( sq_nonneg _ ) ) ( fun n => by rw [ inv_le_inv₀ ] <;> cases abs_cases z <;> nlinarith [ Nat.le_ceil ( |z| ), show ( n : ℝ ) ≥ 0 by positivity ] ) h;
          convert h_summable.mul_left ( 2 ^ 2 / 4 ) using 2 ; push_cast ; ring;
      norm_num [ add_assoc ] at * ; linarith!;
    · contrapose! h;
      refine' summable_nat_add_iff ( ⌈|z|⌉₊ + 1 ) |>.1 _;
      exact Summable.of_nonneg_of_le ( fun n => by positivity ) ( fun n => by rw [ inv_le_comm₀ ] <;> norm_num <;> cases abs_cases z <;> nlinarith [ Nat.le_ceil ( |z| ), show ( n : ℝ ) ≥ 0 by positivity ] ) ( summable_nat_add_iff 1 |>.2 <| Real.summable_one_div_nat_pow.2 one_lt_two )

/-
C expressed in terms of trigamma.
-/
lemma C_sum_eq_trigamma : C_sum = (Real.sqrt 3 / 72) * (trigamma (1/6) + trigamma (1/3) - trigamma (2/3) - trigamma (5/6)) := by
  -- Recognize that the sum of the series can be expressed using the trigamma function.
  have h_trigamma : ∑' n : ℕ, (if n = 0 then 0 else Real.sin (n * Real.pi / 3) / n^2) = (∑' n : ℕ, Real.sin (n * Real.pi / 3) / (n : ℝ)^2) := by
    exact tsum_congr fun n => by aesop;
  -- Recognize that the sum of the series can be expressed using the trigamma function's series definition.
  have h_series : ∑' n : ℕ, Real.sin (n * Real.pi / 3) / (n : ℝ)^2 = (∑' n : ℕ, Real.sin ((6 * n + 1) * Real.pi / 3) / ((6 * n + 1) : ℝ)^2) + (∑' n : ℕ, Real.sin ((6 * n + 2) * Real.pi / 3) / ((6 * n + 2) : ℝ)^2) + (∑' n : ℕ, Real.sin ((6 * n + 4) * Real.pi / 3) / ((6 * n + 4) : ℝ)^2) + (∑' n : ℕ, Real.sin ((6 * n + 5) * Real.pi / 3) / ((6 * n + 5) : ℝ)^2) := by
    have h_split : ∀ {f : ℕ → ℝ}, Summable f → ∑' n : ℕ, f n = (∑' n : ℕ, f (6 * n)) + (∑' n : ℕ, f (6 * n + 1)) + (∑' n : ℕ, f (6 * n + 2)) + (∑' n : ℕ, f (6 * n + 3)) + (∑' n : ℕ, f (6 * n + 4)) + (∑' n : ℕ, f (6 * n + 5)) := by
      intros f hf
      have h_split : ∀ N : ℕ, ∑ n ∈ Finset.range (6 * N), f n = ∑ n ∈ Finset.range N, f (6 * n) + ∑ n ∈ Finset.range N, f (6 * n + 1) + ∑ n ∈ Finset.range N, f (6 * n + 2) + ∑ n ∈ Finset.range N, f (6 * n + 3) + ∑ n ∈ Finset.range N, f (6 * n + 4) + ∑ n ∈ Finset.range N, f (6 * n + 5) := by
        intro N; induction N <;> simp_all +decide [ Nat.mul_succ, Finset.sum_range_succ ] ; ring;
      have h_split : Filter.Tendsto (fun N => ∑ n ∈ Finset.range (6 * N), f n) Filter.atTop (nhds (∑' n : ℕ, f n)) := by
        exact hf.hasSum.tendsto_sum_nat.comp <| Filter.tendsto_id.nsmul_atTop <| by norm_num;
      exact tendsto_nhds_unique h_split ( by simpa only [ * ] using Filter.Tendsto.add ( Filter.Tendsto.add ( Filter.Tendsto.add ( Filter.Tendsto.add ( Filter.Tendsto.add ( hf.comp_injective ( by intros a b; aesop ) |> Summable.hasSum |> HasSum.tendsto_sum_nat ) ( hf.comp_injective ( by intros a b; aesop ) |> Summable.hasSum |> HasSum.tendsto_sum_nat ) ) ( hf.comp_injective ( by intros a b; aesop ) |> Summable.hasSum |> HasSum.tendsto_sum_nat ) ) ( hf.comp_injective ( by intros a b; aesop ) |> Summable.hasSum |> HasSum.tendsto_sum_nat ) ) ( hf.comp_injective ( by intros a b; aesop ) |> Summable.hasSum |> HasSum.tendsto_sum_nat ) ) ( hf.comp_injective ( by intros a b; aesop ) |> Summable.hasSum |> HasSum.tendsto_sum_nat ) );
    convert h_split _ using 1;
    · norm_num [ add_mul, add_div ];
      norm_num [ show ∀ n : ℕ, 6 * n * Real.pi / 3 = 2 * n * Real.pi by intros; ring ];
      norm_num [ Real.sin_add, mul_assoc, mul_left_comm ];
      norm_num [ show ∀ n : ℕ, Real.sin ( n * ( 2 * Real.pi ) ) = 0 from fun n => Real.sin_eq_zero_iff.mpr ⟨ n * 2, by push_cast; ring ⟩ ];
    · -- Since $|\sin(n \pi / 3)| \leq 1$ for all $n$, we have $|\sin(n \pi / 3) / n^2| \leq 1 / n^2$.
      have h_abs : ∀ n : ℕ, |Real.sin (n * Real.pi / 3) / (n : ℝ)^2| ≤ 1 / (n : ℝ)^2 := by
        exact fun n => by simpa [ abs_div ] using div_le_div_of_nonneg_right ( Real.abs_sin_le_one _ ) ( sq_nonneg _ ) ;
      -- Apply the comparison test with the convergent p-series ∑ 1/n².
      have h_comparison : Summable (fun n : ℕ => |Real.sin (n * Real.pi / 3) / (n : ℝ)^2|) := by
        exact Summable.of_nonneg_of_le ( fun n => abs_nonneg _ ) h_abs ( Real.summable_one_div_nat_pow.2 one_lt_two );
      exact h_comparison.of_abs;
  -- Evaluate each of the four series individually.
  have h_series_eval : (∑' n : ℕ, Real.sin ((6 * n + 1) * Real.pi / 3) / ((6 * n + 1) : ℝ)^2) = Real.sqrt 3 / 2 * (∑' n : ℕ, 1 / ((6 * n + 1) : ℝ)^2) ∧
                      (∑' n : ℕ, Real.sin ((6 * n + 2) * Real.pi / 3) / ((6 * n + 2) : ℝ)^2) = Real.sqrt 3 / 2 * (∑' n : ℕ, 1 / ((6 * n + 2) : ℝ)^2) ∧
                      (∑' n : ℕ, Real.sin ((6 * n + 4) * Real.pi / 3) / ((6 * n + 4) : ℝ)^2) = -Real.sqrt 3 / 2 * (∑' n : ℕ, 1 / ((6 * n + 4) : ℝ)^2) ∧
                      (∑' n : ℕ, Real.sin ((6 * n + 5) * Real.pi / 3) / ((6 * n + 5) : ℝ)^2) = -Real.sqrt 3 / 2 * (∑' n : ℕ, 1 / ((6 * n + 5) : ℝ)^2) := by
                        refine' ⟨ _, _, _, _ ⟩ <;> rw [ ← tsum_mul_left ] <;> congr <;> ext n <;> ring;
                        · norm_num [ mul_div, mul_assoc, mul_comm Real.pi, Real.sin_add ] ; ring;
                          norm_num [ mul_two, Real.sin_add, mul_div ] ; ring;
                        · rw [ show Real.pi * ( 2 / 3 ) = Real.pi - Real.pi / 3 by ring ] ; norm_num [ Real.sin_add, mul_two ] ; ring;
                          norm_num [ mul_assoc, mul_comm Real.pi ];
                        · rw [ show Real.pi * ( 4 / 3 ) = Real.pi + Real.pi / 3 by ring, Real.sin_add ] ; norm_num [ mul_two, Real.sin_add, Real.cos_add ] ; ring;
                          norm_num [ Real.cos_sq' ];
                        · norm_num [ ( by ring : Real.pi * ( 5 / 3 ) = 2 * Real.pi - Real.pi / 3 ), Real.sin_add, mul_two ] ; ring;
                          norm_num [ mul_assoc, mul_comm Real.pi ];
  -- Recognize that the sums of the series can be expressed using the trigamma function.
  have h_trigamma_series : (∑' n : ℕ, 1 / ((6 * n + 1) : ℝ)^2) = trigamma (1 / 6) / 36 ∧
                         (∑' n : ℕ, 1 / ((6 * n + 2) : ℝ)^2) = trigamma (1 / 3) / 36 ∧
                         (∑' n : ℕ, 1 / ((6 * n + 4) : ℝ)^2) = trigamma (2 / 3) / 36 ∧
                         (∑' n : ℕ, 1 / ((6 * n + 5) : ℝ)^2) = trigamma (5 / 6) / 36 := by
                           unfold trigamma; norm_num [ div_eq_mul_inv, tsum_mul_left ] ; ring; norm_num;
                           exact ⟨ by rw [ ← tsum_mul_right ] ; exact tsum_congr fun n => by rw [ inv_eq_iff_eq_inv ] ; norm_num ; ring, by rw [ ← tsum_mul_right ] ; exact tsum_congr fun n => by rw [ inv_eq_iff_eq_inv ] ; norm_num ; ring, by rw [ ← tsum_mul_right ] ; exact tsum_congr fun n => by rw [ inv_eq_iff_eq_inv ] ; norm_num ; ring, by rw [ ← tsum_mul_right ] ; exact tsum_congr fun n => by rw [ inv_eq_iff_eq_inv ] ; norm_num ; ring ⟩ ;
  convert h_trigamma using 1 ; rw [ h_series ] ; rw [ h_series_eval.1, h_series_eval.2.1, h_series_eval.2.2.1, h_series_eval.2.2.2 ] ; rw [ h_trigamma_series.1, h_trigamma_series.2.1, h_trigamma_series.2.2.1, h_trigamma_series.2.2.2 ] ; ring;

/-
Simplified value of C.
-/
lemma C_sum_value : C_sum = (Real.sqrt 3 / 48) * (3 * trigamma (1/3) - trigamma (5/6)) := by
  rw [ C_sum_eq_trigamma ];
  -- Applying the duplication formula to trigamma(1/6) and trigamma(1/3), we get:
  have h_duplication : trigamma (1 / 6) + trigamma (2 / 3) = 4 * trigamma (1 / 3) ∧ trigamma (1 / 3) + trigamma (5 / 6) = 4 * trigamma (2 / 3) := by
    constructor <;> have := trigamma_duplication ( 1 / 6 ) <;> have := trigamma_duplication ( 1 / 3 ) <;> have := trigamma_duplication ( 2 / 3 ) <;> norm_num at * <;> linarith!;
  grind +ring

/-
Check riemannZeta type.
-/
#check riemannZeta

/-
Sum of a_n/n^3 is -1/4 * zeta(3) (complex version).
-/
lemma sum_a_n_div_n_pow_3_eq_complex :
  ∑' n : ℕ, (if n = 0 then (0 : ℂ) else (S_a_n n : ℂ) / (n : ℂ)^3) = -1/4 * riemannZeta 3 := by
    -- Apply the formula for `S_a_n` to rewrite the sum.
    have h_sum_formula : ∑' n : ℕ, (if n = 0 then 0 else S_a_n n / (n : ℝ)^3) = ∑' n : ℕ, (-1 / 4 : ℝ) * (if n = 0 then 0 else 1 / (n : ℝ)^3) + ∑' n : ℕ, (-3 / 4 : ℝ) * (if n = 0 then 0 else (if n % 3 = 0 then 1 else 0) / (n : ℝ)^3) + ∑' n : ℕ, (1 : ℝ) * (if n = 0 then 0 else (if n % 4 = 0 then 1 else 0) / (n : ℝ)^3) + ∑' n : ℕ, (3 : ℝ) * (if n = 0 then 0 else (if n % 6 = 0 then 1 else 0) / (n : ℝ)^3) + ∑' n : ℕ, (-3 : ℝ) * (if n = 0 then 0 else (if n % 12 = 0 then 1 else 0) / (n : ℝ)^3) := by
      have h_split_sum : ∀ n : ℕ, (if n = 0 then 0 else S_a_n n / (n : ℝ)^3) = (-1 / 4 : ℝ) * (if n = 0 then 0 else 1 / (n : ℝ)^3) + (-3 / 4 : ℝ) * (if n = 0 then 0 else (if n % 3 = 0 then 1 else 0) / (n : ℝ)^3) + (1 : ℝ) * (if n = 0 then 0 else (if n % 4 = 0 then 1 else 0) / (n : ℝ)^3) + (3 : ℝ) * (if n = 0 then 0 else (if n % 6 = 0 then 1 else 0) / (n : ℝ)^3) + (-3 : ℝ) * (if n = 0 then 0 else (if n % 12 = 0 then 1 else 0) / (n : ℝ)^3) := by
        intro n
        simp [S_a_n_formula];
        split_ifs <;> ring;
      have h_summable : Summable (fun n : ℕ => (-1 / 4 : ℝ) * (if n = 0 then 0 else 1 / (n : ℝ)^3)) ∧ Summable (fun n : ℕ => (-3 / 4 : ℝ) * (if n = 0 then 0 else (if n % 3 = 0 then 1 else 0) / (n : ℝ)^3)) ∧ Summable (fun n : ℕ => (1 : ℝ) * (if n = 0 then 0 else (if n % 4 = 0 then 1 else 0) / (n : ℝ)^3)) ∧ Summable (fun n : ℕ => (3 : ℝ) * (if n = 0 then 0 else (if n % 6 = 0 then 1 else 0) / (n : ℝ)^3)) ∧ Summable (fun n : ℕ => (-3 : ℝ) * (if n = 0 then 0 else (if n % 12 = 0 then 1 else 0) / (n : ℝ)^3)) := by
        refine' ⟨ _, _, _, _, _ ⟩;
        · exact Summable.mul_left _ <| by exact Summable.of_nonneg_of_le ( fun n => by positivity ) ( fun n => by aesop ) <| Real.summable_one_div_nat_pow.2 ( by norm_num : 1 < 3 ) ;
        · refine' Summable.mul_left _ _;
          exact Summable.of_nonneg_of_le ( fun n => by positivity ) ( fun n => by split_ifs <;> norm_num ) ( Real.summable_one_div_nat_pow.2 ( by norm_num : 1 < 3 ) );
        · refine' Summable.of_nonneg_of_le ( fun n => _ ) ( fun n => _ ) ( Real.summable_one_div_nat_pow.2 ( by norm_num : 1 < 3 ) ) <;> aesop;
        · refine' Summable.mul_left _ _;
          exact Summable.of_nonneg_of_le ( fun n => by positivity ) ( fun n => by split_ifs <;> norm_num ) ( Real.summable_one_div_nat_pow.2 ( by norm_num : 1 < 3 ) );
        · refine' Summable.mul_left _ _;
          exact Summable.of_nonneg_of_le ( fun n => by positivity ) ( fun n => by split_ifs <;> norm_num ) ( Real.summable_one_div_nat_pow.2 ( by norm_num : 1 < 3 ) );
      rw [ ← Summable.tsum_add, ← Summable.tsum_add, ← Summable.tsum_add, ← Summable.tsum_add ];
      any_goals tauto;
      · exact tsum_congr h_split_sum;
      · exact Summable.add ( Summable.add ( Summable.add h_summable.1 h_summable.2.1 ) h_summable.2.2.1 ) h_summable.2.2.2.1;
      · exact Summable.add ( Summable.add h_summable.1 h_summable.2.1 ) h_summable.2.2.1;
      · exact Summable.add h_summable.1 h_summable.2.1;
    -- Evaluate each of the sums using the sum_indicator_div_pow lemma.
    have h_sum_eval : (∑' n : ℕ, (-3 / 4 : ℝ) * (if n = 0 then 0 else (if n % 3 = 0 then 1 else 0) / (n : ℝ)^3)) = (-3 / 4 : ℝ) * (3 ^ (-3 : ℂ)) * riemannZeta 3 ∧
                       (∑' n : ℕ, (1 : ℝ) * (if n = 0 then 0 else (if n % 4 = 0 then 1 else 0) / (n : ℝ)^3)) = (1 : ℝ) * (4 ^ (-3 : ℂ)) * riemannZeta 3 ∧
                       (∑' n : ℕ, (3 : ℝ) * (if n = 0 then 0 else (if n % 6 = 0 then 1 else 0) / (n : ℝ)^3)) = (3 : ℝ) * (6 ^ (-3 : ℂ)) * riemannZeta 3 ∧
                       (∑' n : ℕ, (-3 : ℝ) * (if n = 0 then 0 else (if n % 12 = 0 then 1 else 0) / (n : ℝ)^3)) = (-3 : ℝ) * (12 ^ (-3 : ℂ)) * riemannZeta 3 := by
                         refine' ⟨ _, _, _, _ ⟩;
                         · convert congr_arg ( fun x : ℂ => ( -3 / 4 : ℂ ) * x ) ( sum_indicator_div_pow 3 ( by norm_num ) 3 ( by norm_num ) ) using 1 <;> norm_num [ div_eq_mul_inv, tsum_mul_left ] ; ring;
                           · rw [ Complex.ofReal_tsum ] ; norm_num [ mul_assoc, mul_comm, mul_left_comm, ← tsum_mul_left ] ; ring;
                             exact tsum_congr fun n => by split_ifs <;> norm_num;
                           · ring;
                         · convert sum_indicator_div_pow 4 ( by norm_num ) 3 ( by norm_num ) using 1;
                           · norm_num [ Complex.ofReal_tsum ];
                             exact tsum_congr fun n => by split_ifs <;> norm_num;
                           · norm_num;
                         · convert congr_arg ( fun x : ℂ => ( 3 : ℂ ) * x ) ( sum_indicator_div_pow 6 ( by norm_num ) 3 ( by norm_num ) ) using 1 ; norm_num [ div_eq_mul_inv, tsum_mul_left ] ; ring;
                           · rw [ Complex.ofReal_tsum ] ; rw [ ← tsum_mul_right ] ; congr ; ext n ; aesop;
                           · norm_num [ mul_assoc ];
                         · convert congr_arg ( fun x : ℂ => x * ( -3 : ℂ ) ) ( sum_indicator_div_pow 12 ( by norm_num ) 3 ( by norm_num ) ) using 1 <;> norm_num [ div_eq_mul_inv, tsum_mul_left, tsum_mul_right ] ; ring;
                           · rw [ Complex.ofReal_tsum ] ; norm_num [ mul_assoc, mul_comm, mul_left_comm, ← tsum_mul_left ] ; ring;
                             rw [ ← tsum_neg ] ; congr ; ext n ; aesop;
                           · ring;
    convert congr_arg Complex.ofReal h_sum_formula using 1;
    · rw [ Complex.ofReal_tsum ] ; congr ; ext n ; norm_cast ; aesop;
    · norm_num [ tsum_mul_left ] at *;
      rw [ show ( ∑' n : ℕ, if n = 0 then 0 else - ( 1 / 4 * ( n ^ 3 : ℝ ) ⁻¹ ) ) = - ( 1 / 4 ) * ( ∑' n : ℕ, if n = 0 then 0 else ( n ^ 3 : ℝ ) ⁻¹ ) by rw [ ← tsum_mul_left ] ; exact tsum_congr fun n => by split_ifs <;> ring ] ; norm_num [ h_sum_eval ] ; ring;
      rw [ show ( ∑' n : ℕ, if n = 0 then 0 else ( n : ℝ ) ⁻¹ ^ 3 ) = ∑' n : ℕ, ( n : ℝ ) ⁻¹ ^ 3 by rw [ tsum_congr ] ; aesop ] ; norm_num [ Complex.ofReal_tsum ] ; ring;
      norm_cast ; norm_num ; ring;
      norm_num [ Rat.divInt_eq_div ] ; ring;
      norm_num [ zeta_eq_tsum_one_div_nat_cpow ]

/-
Main theorem.
-/
theorem main_theorem : S_sum = 2 * riemannZeta 3 - (Real.pi * Real.sqrt 3 / 24) * trigamma (1/3) + (Real.pi * Real.sqrt 3 / 72) * trigamma (5/6) := by
  rw [ S_eq_final_form, C_sum_value ];
  -- Substitute the expression for the sum of a_n/n^3 into the equation.
  have h_sum : ∑' n : ℕ, (if n = 0 then 0 else S_a_n n / (n : ℂ)^3) = -1 / 4 * riemannZeta 3 := by
    convert sum_a_n_div_n_pow_3_eq_complex using 1;
  convert congr_arg ( fun x : ℂ => - ( 2 * Real.pi / 3 ) * ( Real.sqrt 3 / 48 * ( 3 * trigamma ( 1 / 3 ) - trigamma ( 5 / 6 ) ) ) - 8 * x ) h_sum using 1 ; norm_num ; ring;
  · rw [ Complex.ofReal_tsum ] ; congr ; ext n ; aesop;
  · ring